\ifdefined\pdfoutput\pdfoutput=1\fi
\documentclass[11pt,letterpaper]{article}

\usepackage{geometry}
\geometry{letterpaper, margin=1in}
\usepackage{amsmath,amssymb,amsthm}
\usepackage{mathtools}
\usepackage[hypertexnames=false]{hyperref}
\usepackage{cite}

\theoremstyle{plain}
\newtheorem{theorem}{Theorem}[section]
\newtheorem{proposition}[theorem]{Proposition}
\newtheorem{corollary}[theorem]{Corollary}
\theoremstyle{remark}
\newtheorem{remark}[theorem]{Remark}

\newcommand{\QQ}{\mathbb{Q}}
\newcommand{\ZZ}{\mathbb{Z}}
\newcommand{\PP}{\mathbb{P}}
\newcommand{\Jac}{\operatorname{Jac}}
\newcommand{\Prym}{\operatorname{Prym}}
\newcommand{\Disc}{\operatorname{Disc}}
\newcommand{\Res}{\operatorname{Res}}
\newcommand{\Gal}{\operatorname{Gal}}
\newcommand{\End}{\operatorname{End}}
\newcommand{\Hom}{\operatorname{Hom}}
\newcommand{\Cl}{\operatorname{Cl}}
\newcommand{\F}{\mathbb{F}}
\DeclareMathOperator{\Sha}{\text{III}}

\title{A quartic cover of \texorpdfstring{$37\mathrm{a}1$}{37a1} and its regular \texorpdfstring{$S_4$}{S4}-closure}
\author{%
Haobo Ma\(^{1}\) \and Wenlin Zhang\(^{2,*}\)\\[0.5em]
\small \(^{1}\)AELF PTE LTD.\\
\small \#14-02, Marina Bay Financial Centre Tower 1, 8 Marina Blvd, Singapore 018981, Singapore\\
\small \texttt{auric@aelf.io}\\
\small \(^{2}\)National University of Singapore\\
\small 21 Lower Kent Ridge Road, Singapore 119077, Singapore\\
\small \texttt{e1327962@u.nus.edu}\\
\small \(^{*}\)Corresponding author
}
\date{}

\begin{document}
\maketitle

\begin{abstract}
The conductor-\(37\) elliptic curve \(E\colon\xi^2+\xi=\lambda^3-\lambda\) admits a degree-\(4\) map \(E\to\PP^1_y\) whose branch locus is governed by the irreducible cubic \(c(y)=256y^3+411y^2+165y+32\). We construct the regular \(S_4\)-closure~\(X\) of this cover and establish the isogeny
\[
\Jac(X)\sim_{\QQ}\Jac(X_A)\times E_{\mathrm{res}}^2\times E^3\times Q^3,
\]
where \(X_A\colon w^2=-y(y-1)c(y)\) is a genus-\(2\) curve, \(E_{\mathrm{res}}\) is the resolvent elliptic quotient of conductor~\(1147\), and \(Q=\Prym(Y/E_{\mathrm{res}})\) is an abelian threefold. The splitting field of \(c(y)\) is the unique cubic subextension of the ray class field of \(\QQ(\sqrt{-111})\) of modulus \(\mathfrak p_3^2\); equivalently, it is both the source of a weight-\(1\) cusp form of level~\(999\) and the \(2\)-torsion field of \(\Jac(X_A)\). We determine the bad reduction of \(X_A\) at \(31\) and \(37\), describe the \(2\)-adic and \(3\)-adic cluster pictures of the branch roots, and prove that the wild part of the conductor at \(3\) is exactly \(\delta_3=1\). Thus a single branch cubic links an explicit quartic cover, a cubic ray class field, and the local arithmetic of a genus-\(2\) Jacobian.
\end{abstract}

\noindent\textbf{2020 Mathematics Subject Classification.}
11G10, 11G15, 11R37, 14H30, 14H40.

\noindent\textbf{Keywords.}
Elliptic curves, Galois covers, Prym varieties, Jacobian decompositions, ray class fields, genus-two Jacobians.

\section{Introduction}

When an elliptic curve \(E\) over~\(\QQ\) admits a morphism to~\(\PP^1\) of small degree, the Galois closure of the corresponding function field extension carries a rich arithmetic structure. The quotient curves of this closure, together with the Kani--Rosen idempotent decomposition of its Jacobian \cite{KaniRosen1989}, yield a family of abelian varieties whose conductors, endomorphism algebras, and \(L\)-functions are controlled by the branch data. In practice, however, making every piece of this picture explicit requires computing the regular closure, identifying each quotient, and matching each isotypic factor to a concrete abelian variety --- a programme that quickly becomes unwieldy as the degree grows.

In this paper we carry out this programme for the degree-\(4\) map defined by
\begin{equation}\label{eq:quartic}
\Pi(\lambda,y)=\lambda^4-\lambda^3-(2y+1)\lambda^2+\lambda+y(y+1)
\end{equation}
and the cubic
\begin{equation}\label{eq:branch-cubic}
c(y)=256y^3+411y^2+165y+32.
\end{equation}
After the substitution \(y=\lambda^2+\xi\), the equation \(\Pi(\lambda,y)=0\) becomes the Weierstrass model \(\xi^2+\xi=\lambda^3-\lambda\) of the conductor-\(37\) elliptic curve (Cremona label \(37\mathrm{a}1\)), and the function \(y=\lambda^2+\xi\) gives a degree-\(4\) cover \(E\to\PP^1_y\) whose finite branch locus is \(\{0,1\}\cup\{c(y)=0\}\).

What makes this example particularly interesting is that the single cubic \(c(y)\) simultaneously governs several apparently unrelated arithmetic objects: its splitting field yields a weight-\(1\) modular form via class field theory, the hyperelliptic curve \(w^2=-y(y-1)c(y)\) contributes a genus-\(2\) Jacobian to the regular closure, and the cubic resolvent of the quartic produces a second elliptic curve. The constructions involved --- quartic resolvents, Jacobian decompositions, Prym varieties, ray class fields --- are all standard (see \cite{Cox2012,KaniRosen1989,Mumford1974,BirkenhakeLange2004,Neukirch1999,CaroccaRodriguez2006,Paulhus2008} for general background), but this cover has the unusual feature that the quotient geometry, the class-field-theoretic cubic, the weight-\(1\) Artin form, and the local branch-root data can all be tracked explicitly from the same branch polynomial.

The contribution of the paper is therefore explicit rather than abstract. The works of Kani--Rosen \cite{KaniRosen1989} and Carocca--Rodr{\'i}guez \cite{CaroccaRodriguez2006} provide the rational-idempotent formalism for Jacobians with group actions, Paulhus \cite{Paulhus2008} studies decomposition patterns for curves with extra automorphisms, and Mumford's and Lange--Ortega's work \cite{Mumford1974,LangeOrtega2011} develops the geometric theory of Prym varieties. What is new here is that one regular \(S_4\)-quartic cover can be carried through all of these frameworks in closed form, and that the resulting package is large enough to exhibit several phenomena simultaneously: a nontrivial genus-\(2\) factor, a split Prym surface, a cubic ray class field and weight-\(1\) form, and a wild prime whose contribution can still be isolated exactly. In that sense the example is meant to be informative rather than anecdotal.

The paper should therefore be read as a fully worked arithmetic example with a transferable method. Starting from one explicit quartic cover, we recover the regular \(S_4\)-closure and its quotients, apply the rational-idempotent method to decompose the Jacobian, identify the cubic splitting field by class-field-theoretic means, and then feed the \(p\)-adic root data into the conductor formulas for the genus-\(2\) quotient. The point is not a new abstract decomposition theorem, but a complete template for the class of regular \(S_4\)-quartic covers whose branch cubic and resolvent geometry remain explicit.

The logical status of the results is as follows. The structural assertions of the paper --- the \(S_4\)-closure, the quotient curves, the Jacobian decomposition, the class-field-theoretic description of the splitting field, and the exact wild contribution \(\delta_3=1\) at the prime \(3\) --- are proved in the text from standard geometric, representation-theoretic, and local-arithmetic input. The companion scripts are used only for finite-field point counts, Frobenius tables, and independent audit files for the explicit calculations recorded in Sections~5 and~6.

\subsection{Summary of results}
We show that \(\Pi(\lambda,y)\) has Galois group \(S_4\) over \(\QQ(y)\), and that the regular closure \(X\) has genus \(16\) (Theorem~\ref{thm:closure}). The relevant quotients are the genus-\(2\) curve \(X_A=X/A_4\) given by \(w^2=-y(y-1)c(y)\), the genus-\(4\) curve \(X_V=X/V_4\), and the resolvent elliptic curve \(E_{\mathrm{res}}=X/D_4\), which turns out to be birational to the curve \(w_{\mathrm{res}}^2=x_{\mathrm{res}}^3-16x_{\mathrm{res}}^2-64x_{\mathrm{res}}+1040\) of conductor \(1147=31\cdot 37\) (Proposition~\ref{prop:eres}). The map \(X_V\to X_A\) is an \'etale cyclic triple cover (Theorem~\ref{thm:quotients}).

Using the \(S_4\)-representation structure on \(H^0(X,\Omega_X^1)\), we establish the isogeny
\[
\Jac(X)\sim_{\QQ}\Jac(X_A)\times E_{\mathrm{res}}^2\times E^3\times Q^3,
\]
where \(Q=\Prym(Y/E_{\mathrm{res}})\) is an abelian threefold with \(Y=X/C_4\) (Theorem~\ref{thm:full-factor}). The Prym surface \(P=\Prym(X_V/X_A)\) satisfies \(P\sim_{\QQ}E_{\mathrm{res}}^2\) and has endomorphism algebra \(M_2(\QQ)\) (Proposition~\ref{prop:prym-end}).

The genus-\(2\) factor turns out to be the most delicate part of the decomposition. We record the Frobenius characteristic polynomials of \(\Jac(X_A)\) at the good primes \(5\), \(13\), and \(17\), and use additional good primes up to \(53\) as irreducibility witnesses for the local quartic factors (Proposition~\ref{thm:simple}). These data provide a concrete arithmetic profile for the genus-\(2\) factor and contrast sharply with the explicitly split Prym surface \(E_{\mathrm{res}}^2\).

On the arithmetic side, we show that the splitting field \(L\) of \(c(y)\) sits inside the ray class field of \(K=\QQ(\sqrt{-111})\) of modulus \(\mathfrak{p}_3^2\), where \(\mathfrak{p}_3\) is the unique prime above~\(3\) (Theorem~\ref{thm:ray-class}). The ray class group \(\Cl_{\mathfrak{p}_3^2}(K)\) is cyclic of order~\(24\), and \(L\) is its unique cubic subextension. The Artin representation of \(L/\QQ\) gives a weight-\(1\) cuspidal newform of level \(999=3^3\cdot 37\), whose Hecke eigenvalues at unramified primes are \(a_p=\#\{y\in\F_p:c(y)=0\}-1\) (Theorem~\ref{thm:hecke-traces}).

We further show that the elliptic quotients \(E\) and \(E_{\mathrm{res}}\) are not \(\QQ\)-isogenous and that the product decomposition forces \(\rho(\Jac(X)_{\overline{\QQ}})\geq 16\) (Proposition~\ref{thm:non-isogeny} and Corollary~\ref{cor:picard}). At the two tame primes of bad reduction for~\(X_A\), the nodal singularities have opposite splitting types: the node at~\(31\) is split while the node at~\(37\) is non-split (Proposition~\ref{prop:node-type}). The \(2\)-torsion field \(\QQ(\Jac(X_A)[2])\) coincides with~\(L\), and the mod-\(2\) Galois image is the small group \(S_3\subset\mathrm{Sp}(4,\F_2)\cong S_6\) of order~\(6\) (Theorem~\ref{thm:two-torsion}).

We also determine the \(\F_2[S_3]\)-module structure of the \(2\)-torsion: the Weil pairing admits a Galois-invariant orthogonal splitting into two hyperbolic planes, one fixed and one carrying the standard representation of~\(S_3\cong\mathrm{GL}(2,\F_2)\) (Theorem~\ref{thm:mod2-structure}). At the semistable primes, the component groups and Tamagawa numbers of \(\Jac(X_A)\) match those of \(E_{\mathrm{res}}\) exactly (Proposition~\ref{prop:tamagawa}).

At the two wild primes, the \(p\)-adic cluster pictures of the branch roots are determined explicitly (Proposition~\ref{prop:cluster-pictures}). At \(p=3\), the three roots of \(c(y)\) coalesce to a single cluster at fractional depth~\(3/2\), requiring a totally ramified cubic extension of~\(\QQ_3\) to resolve; the hyperelliptic conductor formula then gives the exact wild contribution \(\delta_3=1\) to the conductor at~\(3\) (Theorem~\ref{thm:conductor-exact}). At \(p=2\), two of the three roots pair with the rational roots \(0\) and~\(1\) in balanced clusters of depth~\(5\), while the third root is isolated at valuation~\(-8\).

\subsection{Organization}
In Section~2 we recover the elliptic model of the quartic and study the attached cubic field. Section~3 constructs the regular \(S_4\)-closure and its quotient curves. In Section~4 we establish the Jacobian decomposition via representation-theoretic methods and determine the endomorphism algebra of the Prym surface. Section~5 is devoted to the ray class field and the associated weight-\(1\) modular form. In Section~6 we study the arithmetic of the genus-\(2\) factor: Frobenius data, automorphic structure, Picard-number bounds, local conductor calculations, the \(\F_2\)-module structure of the \(2\)-torsion, and numerical verification. Section~7 records several further directions suggested by the computation.

\section{The quartic cover}

In this section we recover the elliptic model of the quartic and extract the cubic field attached to the finite branch locus.

\subsection{The elliptic model}

\begin{theorem}\label{thm:elliptic}
Under the change of variables
\[
\xi=y-\lambda^2,
\]
the affine curve \(\Pi(\lambda,y)=0\) becomes
\begin{equation}\label{eq:elliptic}
E:\quad \xi^2+\xi=\lambda^3-\lambda.
\end{equation}
Its smooth projective model is an elliptic curve over \(\QQ\) with minimal discriminant \(37\), \(j\)-invariant \(110592/37\), and Cremona label \(37\mathrm{a}1\).
\end{theorem}

\begin{proof}
Treating \(\Pi\) as a quadratic in \(y\) gives
\[
y^2+(1-2\lambda^2)y+\lambda^4-\lambda^3-\lambda^2+\lambda=0.
\]
Substituting \(y=\lambda^2+\xi\) gives
\[
\xi^2+\xi=\lambda^3-\lambda.
\]
This is a minimal Weierstrass model with the stated invariants \cite{Silverman2009,LMFDBEllipticCurve37a1}.
\end{proof}

The rational function
\begin{equation}\label{eq:y-function}
y=\lambda^2+\xi\in \QQ(E)
\end{equation}
defines a morphism
\[
\varphi:E\longrightarrow \PP^1_y.
\]

\begin{proposition}\label{prop:ramification}
The map \(\varphi\) has degree \(4\). It is totally ramified of index \(4\) above \(y=\infty\), and has exactly five finite simple ramification points. The finite branch locus is
\[
\{0,1\}\cup\{c(y)=0\},
\]
where \(c(y)\) is given by \eqref{eq:branch-cubic}. Equivalently,
\[
\Disc_\lambda \Pi(\lambda,y)=-y(y-1)c(y).
\]
\end{proposition}

\begin{proof}
At infinity, \(\operatorname{ord}_O(\lambda)=-2\) and \(\operatorname{ord}_O(\xi)=-3\), so \(\operatorname{ord}_O(y)=-4\); hence \(\deg\varphi=4\) and \(y=\infty\) is totally ramified.

Differentiating \eqref{eq:elliptic} gives \((2\xi+1)\,d\xi=(3\lambda^2-1)\,d\lambda\), whence
\[
dy=\frac{2\lambda(2\xi+1)+3\lambda^2-1}{2\xi+1}\,d\lambda.
\]
The ramification condition \(dy=0\) requires \(2\lambda(2\xi+1)+3\lambda^2-1=0\), that is, \(\xi=-(3\lambda^2+2\lambda-1)/(4\lambda)\). Substituting into \(\xi^2+\xi=\lambda^3-\lambda\) and clearing denominators gives a degree-\(5\) polynomial in \(\lambda\) that factors as
\[
(\lambda-1)(\lambda+1)(16\lambda^3-9\lambda^2+1)=0.
\]
The roots \(\lambda=\pm 1\) give \(y=0,1\). Eliminating \(\lambda\) between \(16\lambda^3-9\lambda^2+1=0\) and
\[
4\lambda y-(4\lambda^3-3\lambda^2-2\lambda+1)=0
\]
gives
\[
\Res_\lambda\!\bigl(16\lambda^3-9\lambda^2+1,\;4\lambda y-(4\lambda^3-3\lambda^2-2\lambda+1)\bigr)=-64\,c(y).
\]
The discriminant identity is the quartic reformulation.
\end{proof}

\subsection{The cubic field}

\begin{proposition}\label{prop:number-field}
Let \(\beta\) be a root of \(c(y)\), and let \(F=\QQ(\beta)\).
Then:
\begin{enumerate}
\item[(i)] \(\Disc(c)=-3^9\cdot 31^2\cdot 37\).
\item[(ii)] Let \(\alpha\) be a root of
\[
16x^3-9x^2+1,
\]
chosen so that
\[
\beta=\frac{4\alpha^3-3\alpha^2-2\alpha+1}{4\alpha}.
\]
If \(u:=16\alpha\), then \(F=\QQ(u)\), the element \(u\) satisfies
\[
f_u(u)=0,\qquad f_u(t)=t^3-9t^2+256,
\]
and
\[
\theta:=256\beta=-3u^2+24u-128.
\]
\item[(iii)] The field discriminant of \(F\) is
\[
\Disc(F/\QQ)=-999=-3^3\cdot 37.
\]
\item[(iv)] The element
\[
\omega:=\frac{u^2+7u+16}{32}
\]
is integral, and
\[
\mathcal O_F=\ZZ\cdot 1\oplus \ZZ u\oplus \ZZ\omega.
\]
\item[(v)] The integral model \(\theta=256\beta\) satisfies the monic polynomial
\[
g(\theta)=\theta^3+411\theta^2+42240\theta+2097152,
\]
and the index of the monogenic order in \(\mathcal O_F\) is
\[
[\mathcal O_F:\ZZ[\theta]]=2^8\cdot 3^3\cdot 31=214272.
\]
\item[(vi)] The quadratic resolvent field of the splitting field of \(c(y)\) is
\[
K_\Delta=\QQ(\sqrt{-111}).
\]
\end{enumerate}
\end{proposition}

\begin{proof}
Part~(i) is a direct computation from \eqref{eq:branch-cubic}.

For~(ii), the nontrivial ramification points satisfy
\[
16\lambda^3-9\lambda^2+1=0,
\qquad
y=\frac{4\lambda^3-3\lambda^2-2\lambda+1}{4\lambda}.
\]
Choosing a compatible pair \((\alpha,\beta)\) gives \(F=\QQ(\beta)=\QQ(\alpha)\). The scaling \(u=16\alpha\) yields \(f_u(u)=0\), and substituting back gives
\[
\theta=256\beta=u^2-12u-128+\frac{1024}{u}.
\]
Using \(u^3=9u^2-256\), hence \(1024/u=-4u^2+36u\), this simplifies to
\[
\theta=-3u^2+24u-128.
\]

For~(iii)--(iv), one checks that \(\omega=(u^2+7u+16)/32\) satisfies
\[
\omega^3-6\omega^2+15\omega-2=0,
\]
so \(\omega\) is integral. Set
\[
A:=\ZZ\oplus \ZZ u\oplus \ZZ\omega.
\]
The change-of-basis matrix from \((1,u,u^2)\) to \((1,u,\omega)\) has determinant \(1/32\), hence
\[
\Disc(A)=\frac{\Disc(f_u)}{32^2}.
\]
Since \(\Disc(f_u)=-2^{10}\cdot 3^3\cdot 37\), this gives
\[
\Disc(A)=-3^3\cdot 37=-999.
\]

To verify maximality of \(A\), we check locally at \(3\) and~\(37\). Since \(v_{37}(\Disc(A))=1\), we have \(37\nmid [\mathcal O_F:A]\) from
\[
\Disc(A)=\Disc(F/\QQ)\,[\mathcal O_F:A]^2.
\]
At \(3\), set
\[
\alpha:=u-2\omega-1.
\]
From \(u^2+7u+16=32\omega\) and \(u^3-9u^2+256=0\),
\[
\alpha^3+6\alpha^2+3\alpha+6=0,
\]
and
\[
u=-\alpha^2-4\alpha+5,
\qquad
\omega=2-\frac{\alpha^2+5\alpha}{2}.
\]
Since \(2\in\ZZ_3^\times\), it follows that
\[
A\otimes \ZZ_3=\ZZ_3[\alpha].
\]
The polynomial \(x^3+6x^2+3x+6\) is \(3\)-Eisenstein, so \(\ZZ_3[\alpha]\) is the maximal order at~\(3\). Therefore
\[
\mathcal O_F=A=\ZZ\oplus \ZZ u\oplus \ZZ\omega.
\]
Consequently,
\[
\Disc(F/\QQ)=\Disc(A)=-999.
\]

Note for later use that \(3\) is totally ramified in \(F\) (by the Eisenstein property) and \(37\) is tamely ramified (\(v_{37}(\Disc(F))=1\)).

For~(v), \(\theta=256\beta\) satisfies \(g(\theta)=0\), and
\[
\Disc(g)=256^2\,\Disc(c)=-2^{16}\cdot 3^9\cdot 31^2\cdot 37.
\]
Since \(\theta\in \ZZ[u]\), comparing discriminants of the monogenic orders gives
\[
\Disc(\ZZ[u])=\Disc(f_u)=-2^{10}\cdot 3^3\cdot 37
\]
and therefore
\[
[\ZZ[u]:\ZZ[\theta]]
=\sqrt{\frac{\Disc(g)}{\Disc(f_u)}}
=2^3\cdot 3^3\cdot 31
=6696.
\]
Since \(\omega=(u^2+7u+16)/32\), expressing \((1,u,\omega)\) in the basis \((1,u,u^2)\) gives a lower-triangular matrix with diagonal entries \((1,1,1/32)\); its determinant is \(1/32\), whence \([\mathcal O_F:\ZZ[u]]=32\). Therefore
\[
[\mathcal O_F:\ZZ[\theta]]
=[\mathcal O_F:\ZZ[u]]\,[\ZZ[u]:\ZZ[\theta]]
=32\cdot 6696
=2^8\cdot 3^3\cdot 31.
\]

Finally, \(c(y)\) is irreducible with nonsquare discriminant, so its splitting field has Galois group \(S_3\), and its unique quadratic subfield is \(\QQ(\sqrt{\Disc(F)})=\QQ(\sqrt{-999})=\QQ(\sqrt{-111})\).
\end{proof}

\begin{remark}\label{rem:cubic-field}
The field \(F\) is recorded in the LMFDB \cite{LMFDB} as \texttt{3.1.999.1}. Its class number is~\(1\), so the ring of integers \(\mathcal O_F\) is a principal ideal domain. The index
\[
[\mathcal O_F:\ZZ[\theta]]=2^8\cdot 3^3\cdot 31
\]
shows that the natural generator \(\theta=256\beta\) is far from integral; the discrepancy comes from the leading coefficient \(256=2^8\), the ramified prime~\(3\), and the additional prime~\(31\).
\end{remark}

\section{The regular \texorpdfstring{$S_4$}{S4}-closure}

We now determine the Galois closure of the quartic cover and identify the quotient curves that will enter the Jacobian decomposition.

\subsection{Monodromy and genus}

\begin{proposition}\label{prop:s4}
The quartic \(\Pi(\lambda,y)\) has Galois group
\[
\Gal(\Pi/\QQ(y))\cong S_4.
\]
\end{proposition}

\begin{proof}
Since the discriminant \(-y(y-1)c(y)\) is not a square in \(\QQ(y)\), the Galois group is not contained in \(A_4\). Specializing at \(y=2\) gives
\[
\Pi(\lambda,2)=\lambda^4-\lambda^3-5\lambda^2+\lambda+6.
\]
Modulo \(5\) this is irreducible (a \(4\)-cycle), and modulo \(3\) it factors as a linear times an irreducible cubic (a \(3\)-cycle). A transitive subgroup of \(S_4\) containing both cycle types is all of~\(S_4\).
\end{proof}

\begin{theorem}\label{thm:closure}
Let \(X\) be the smooth projective curve with function field equal to the Galois closure of \(\QQ(E)/\QQ(y)\).
Then \(X\to\PP^1_y\) is a regular \(S_4\)-cover of degree \(24\), branched with inertia orders
\[
4,2,2,2,2,2.
\]
In particular,
\[
g(X)=16.
\]
\end{theorem}

\begin{proof}
The cover \(E\to\PP^1_y\) has local monodromy of order~\(4\) at infinity and five transpositions at finite branch points (Proposition~\ref{prop:ramification}). Since the Galois group is \(S_4\) (Proposition~\ref{prop:s4}), the regular closure has degree~\(24\). The inertia groups in the regular cover are generated by the corresponding monodromy elements.

The Hurwitz formula gives
\[
2g(X)-2
=24\left(-2+\left(1-\frac14\right)+5\left(1-\frac12\right)\right)
=30,
\]
hence \(g(X)=16\).
\end{proof}

\subsection{Resolvent quotients}

The cubic resolvent of \eqref{eq:quartic}, built from the three pairings \(z_{ij}=\lambda_i\lambda_j+\lambda_k\lambda_l\) of the four roots \cite{Cox2012}, is
\begin{equation}\label{eq:resolvent}
R(z,y)=z^3+(2y+1)z^2-(4y^2+4y+1)z-(8y^3+13y^2+5y+1).
\end{equation}

\begin{proposition}\label{prop:resolvent}
The discriminant of \(R(z,y)\) with respect to \(z\) equals
\[
\Disc_z R(z,y)=-y(y-1)c(y).
\]
The splitting field of \(R\) over \(\QQ(y)\) is the fixed field of the Klein four subgroup \(V_4\subset S_4\).
\end{proposition}

\begin{proof}
A direct computation gives
\[
\Disc_z R(z,y)=-y(y-1)\bigl(256y^3+411y^2+165y+32\bigr).
\]
The fixed field statement is classical \cite{Cox2012}.
\end{proof}

\begin{theorem}\label{thm:quotients}
Let
\[
X_A:=X/A_4,\qquad X_V:=X/V_4.
\]
Then:
\begin{enumerate}
\item[(i)] \(X_A\) is the genus-\(2\) curve
\[
X_A:\quad w^2=-y(y-1)c(y).
\]
\item[(ii)] \(X_V\) is the genus-\(4\) curve with function field
\[
\QQ(X_V)=\QQ(y,z,\sqrt{-y(y-1)c(y)}),\qquad R(z,y)=0.
\]
\item[(iii)] The natural map \(X_V\to X_A\) is an \'etale cyclic cubic cover.
\end{enumerate}
\end{theorem}

\begin{proof}
The \(A_4\)-quotient is the unique quadratic subextension, generated by the square root of the discriminant:
\[
\QQ(X_A)=\QQ\bigl(y,\sqrt{-y(y-1)c(y)}\bigr).
\]
The right-hand side is squarefree of degree~\(5\), giving genus~\(2\). The \(V_4\)-quotient is the resolvent splitting field (Proposition~\ref{prop:resolvent}).

At each finite branch point the inertia is a transposition, meeting both \(A_4\) and \(V_4\) trivially. At infinity the inertia is generated by a \(4\)-cycle \(\sigma\), and
\[
\langle \sigma\rangle\cap A_4=\langle \sigma^2\rangle=\langle \sigma\rangle\cap V_4.
\]
The relative inertia in \(A_4/V_4\cong C_3\) is therefore trivial everywhere, so \(X_V\to X_A\) is \'etale. Hurwitz gives
\[
2g(X_V)-2=3(2g(X_A)-2)=6,
\]
hence \(g(X_V)=4\).
\end{proof}

\begin{corollary}\label{cor:prym}
The connected norm kernel
\[
P:=\Prym(X_V/X_A)
\]
is an abelian surface. Moreover,
\[
\QQ(\zeta_3)\hookrightarrow \End^0(P),
\qquad
\Jac(X_V)\sim_{\QQ}\Jac(X_A)\times P.
\]
\end{corollary}

\begin{proof}
The cover \(X_V\to X_A\) is \'etale of degree~\(3\) with \(g(X_V)=4\), so \(\dim P=4-2=2\). The standard Prym decomposition for a cyclic triple cover gives
\[
\Jac(X_V)\sim_{\QQ}\Jac(X_A)\times P
\]
and the deck action furnishes \(\QQ(\zeta_3)\hookrightarrow \End^0(P)\) \cite[Section~12.9]{BirkenhakeLange2004}.
\end{proof}

\section{The Jacobian}

Our goal is to decompose \(\Jac(X)\) up to \(\QQ\)-isogeny by exploiting the \(S_4\)-module structure on \(H^0(X,\Omega_X^1)\). We first work out this structure and then identify each isotypic factor with a concrete abelian variety.

\subsection{Differentials}

\begin{theorem}\label{thm:hodge}
Let \(V_2\) be the two-dimensional irreducible representation of \(S_4\), \(V_3\) the standard three-dimensional representation, and \(V_3^-:=V_3\otimes\mathrm{sgn}\).
Then
\[
H^0(X,\Omega_X^1)\cong \mathrm{sgn}^{\oplus 2}\oplus V_2\oplus V_3\oplus (V_3^-)^{\oplus 3}.
\]
Accordingly there are rational isotypic factors \(A_{\mathrm{sgn}}, A_{V_2}, A_{V_3}\), and \(A_{V_3^-}\) such that
\[
\Jac(X)\sim_{\QQ}A_{\mathrm{sgn}}\times A_{V_2}\times A_{V_3}\times A_{V_3^-},
\]
with dimensions
\[
\dim A_{\mathrm{sgn}}=2,\qquad \dim A_{V_2}=2,\qquad \dim A_{V_3}=3,\qquad \dim A_{V_3^-}=9.
\]
\end{theorem}

\begin{proof}
Since \(X/S_4\cong \PP^1\), there is no trivial summand. Write
\[
H^0(X,\Omega_X^1)\cong m_{\mathrm{sgn}}\cdot \mathrm{sgn}\oplus m_2\cdot V_2\oplus m_3\cdot V_3\oplus m_3^-\cdot V_3^-.
\]
For any subgroup \(H\subset S_4\), the \(H\)-invariants have dimension \(g(X/H)\).

From Theorem~\ref{thm:quotients}, \(g(X/A_4)=2\) and \(g(X/V_4)=4\). If \(S_3\subset S_4\) is a point stabilizer, then \(X/S_3\cong E\), so \(g(X/S_3)=1\). Let \(C=\langle(1234)\rangle\) and \(D=N_{S_4}(C)\cong D_4\). Character averages give
\[
\begin{array}{c|ccccc}
 & A_4 & V_4 & S_3 & C & D\\ \hline
\mathrm{sgn} & 1 & 1 & 0 & 0 & 0\\
V_2 & 0 & 2 & 0 & 1 & 1\\
V_3 & 0 & 0 & 1 & 0 & 0\\
V_3^- & 0 & 0 & 0 & 1 & 0
\end{array}
\]
for the invariant dimensions of each irreducible constituent.

Reading off the \(A_4\)- and \(V_4\)-columns gives
\[
m_{\mathrm{sgn}}=2,\qquad m_{\mathrm{sgn}}+2m_2=4,
\]
hence \(m_2=1\), and the \(S_3\)-column gives \(m_3=1\). The genus constraint \(g(X)=16\) then forces
\[
16=m_{\mathrm{sgn}}+2m_2+3m_3+3m_3^-,
\]
so \(m_3^-=3\). The Jacobian decomposition follows from the rational isotypic decomposition.
\end{proof}

\subsection{The dihedral quotient}

\begin{proposition}\label{prop:eres}
Let \(C=\langle(1234)\rangle\subset S_4\), let \(D=N_{S_4}(C)\cong D_4\), and set
\[
Y:=X/C,\qquad E_{\mathrm{res}}:=X/D.
\]
Then:
\begin{enumerate}
\item[(i)] The function field of \(E_{\mathrm{res}}\) is
\[
\QQ(E_{\mathrm{res}})=\QQ(y,z),\qquad R(z,y)=0,
\]
where \(R\) is the cubic resolvent \eqref{eq:resolvent}.
\item[(ii)] The smooth projective model of \(R(z,y)=0\) is birational over \(\QQ\) to the elliptic curve
\[
 \mathcal E_{\mathrm{res}}:\quad w_{\mathrm{res}}^2=x_{\mathrm{res}}^3-16x_{\mathrm{res}}^2-64x_{\mathrm{res}}+1040.
\]
In particular,
\[
g(E_{\mathrm{res}})=1,\qquad
\Delta(\mathcal E_{\mathrm{res}})=2^{12}\cdot 31^2\cdot 37,\qquad
j(\mathcal E_{\mathrm{res}})=\frac{2^{18}\cdot 7^3}{31^2\cdot 37}.
\]
\item[(iii)] \(g(Y)=4\), and the natural map \(Y\to E_{\mathrm{res}}\) has degree \(2\) and is branched at exactly six points.
\end{enumerate}
\end{proposition}

\begin{proof}
Since \(V_4\triangleleft D\) and \(D/V_4\) has order \(2\), the quotient \(E_{\mathrm{res}}=X/D\) sits inside the resolvent field \(\QQ(X_V)/\QQ(y)\). Under the identification \(S_4/V_4\cong S_3\), the image of \(D\) stabilizes one of the three pairings of the four quartic roots, so the fixed field is generated by a single resolvent root \(z\):
\[
\QQ(E_{\mathrm{res}})=\QQ(y,z),\qquad R(z,y)=0.
\]

The affine model \(R(z,y)=0\) passes through \((y,z)=(0,1)\). Projecting from this point, set
\[
u=\frac{z-1}{y}.
\]
Substituting \(z=1+uy\) into \(R(z,y)=0\) gives
\[
(u^3+2u^2-4u-8)y^2+(4u^2-17)y+(4u-7)=0.
\]
Set
\[
v=2(u^3+2u^2-4u-8)y+(4u^2-17).
\]
Then the quadratic discriminant identity becomes
\[
v^2=-4u^3-16u^2+16u+65.
\]
After the change of variables
\[
 x_{\mathrm{res}}=-4u,\qquad w_{\mathrm{res}}=4v,
\]
we obtain the Weierstrass model
\[
 w_{\mathrm{res}}^2=x_{\mathrm{res}}^3-16x_{\mathrm{res}}^2-64x_{\mathrm{res}}+1040.
\]
The inverse rational map is recovered from
\[
 u=-\frac{x_{\mathrm{res}}}{4},\qquad
 v=\frac{w_{\mathrm{res}}}{4},\qquad
 y=\frac{v-(4u^2-17)}{2(u^3+2u^2-4u-8)},\qquad
 z=1+uy.
\]
The normalization of \(R(z,y)=0\) is therefore an elliptic curve, and the discriminant and \(j\)-invariant are read off from the integral model.

From the character table in Theorem~\ref{thm:hodge}, the invariant dimensions for \(C\) and \(D\) are
\[
g(Y)=\dim H^0(X,\Omega_X^1)^C=1+3=4,\qquad
g(E_{\mathrm{res}})=\dim H^0(X,\Omega_X^1)^D=1.
\]
Since \([D:C]=2\), the map \(Y\to E_{\mathrm{res}}\) has degree \(2\). Hurwitz gives
\[
2g(Y)-2=2(2g(E_{\mathrm{res}})-2)+r=r,
\]
so the total ramification degree is \(r=6\).
\end{proof}

\begin{proposition}\label{prop:eres-conductor}
The minimal discriminant and conductor of \(\mathcal E_{\mathrm{res}}\) are
\[
\Delta_{\min}(\mathcal E_{\mathrm{res}})=31^2\cdot 37,
\qquad
N(\mathcal E_{\mathrm{res}})=31\cdot 37=1147.
\]
In particular, \(\mathcal E_{\mathrm{res}}\) has multiplicative reduction at both \(31\) and \(37\).
\end{proposition}

\begin{proof}
A short computation gives \(c_4=2^{10}\cdot 7\), \(c_6=-2^9\cdot 23\cdot 29\), and \(\Delta=2^{12}\cdot 31^2\cdot 37\).
Since \(v_2(c_4)=10\ge 4\), \(v_2(c_6)=9\ge 6\), and \(v_2(\Delta)=12\ge 12\), one Tate reduction step with \(u=2\) yields the minimal model with
\[
c_4'=c_4/2^4=448=2^6\cdot 7,
\qquad
c_6'=c_6/2^6=-5336=-2^3\cdot 23\cdot 29,
\qquad
\Delta'=\Delta/2^{12}=31^2\cdot 37.
\]
Since \(v_2(c_6')=3<6\), no further reduction at~\(2\) is possible; at all odd primes, \(v_p(c_4')<4\) or \(v_p(\Delta')=0\), so the model is globally minimal.
Now \(v_{31}(c_4')=0\) and \(v_{31}(\Delta')=2>0\), so the reduction at \(31\) is multiplicative (conductor exponent~\(1\)).
Similarly \(v_{37}(c_4')=0\) and \(v_{37}(\Delta')=1>0\), giving multiplicative reduction at \(37\).
\end{proof}

\begin{proposition}\label{prop:extract}
Let \(G\) be a finite group acting on a smooth projective curve \(X\) over \(\QQ\), and let \(W\) be an absolutely irreducible rational representation of \(G\) such that \(\End_G(W)=\QQ\). Write the \(W\)-isotypic component of \(H^0(X,\Omega_X^1)\) as
\[
H^0(X,\Omega_X^1)_W\cong W\otimes M_W,
\]
where \(M_W\) is the multiplicity space.
If \(A_W\subset \Jac(X)\) denotes the corresponding rational isotypic factor, then there exists an abelian variety \(B_W\), unique up to \(\QQ\)-isogeny, such that
\[
H^0(B_W,\Omega_{B_W}^1)\cong M_W,
\qquad
A_W\sim_{\QQ} B_W^{\dim W}.
\]
In particular, whenever \(M_W\) is realized by the holomorphic differentials of a quotient or Prym variety \(B\), one has
\[
A_W\sim_{\QQ} B^{\dim W}.
\]
\end{proposition}

\begin{proof}
We recall the standard group-algebra argument underlying \cite{KaniRosen1989}; see also \cite[Section~13.4]{BirkenhakeLange2004}.

Let \(e_W\in \QQ[G]\) be the central idempotent attached to \(W\). Since \(\End_G(W)=\QQ\), the corresponding Wedderburn block is
\[
\QQ[G]e_W\cong M_d(\QQ),\qquad d=\dim W.
\]
Choose a primitive idempotent \(\varepsilon\in M_d(\QQ)\), and define
\[
B_W:=\varepsilon\cdot A_W.
\]
The central idempotent \(e_W\) is a sum of \(d\) pairwise orthogonal conjugates of \(\varepsilon\), so
\[
A_W\sim_{\QQ} B_W^d.
\]
On differentials, \(H^0(A_W,\Omega_{A_W}^1)\cong W\otimes M_W\), and \(\varepsilon\) picks out a single matrix-column factor:
\[
H^0(B_W,\Omega_{B_W}^1)\cong M_W.
\]
Any other abelian variety \(B\) realizing the same multiplicity space defines the same rational Hodge substructure of \(H^1(X)\), hence is \(\QQ\)-isogenous to~\(B_W\).
\end{proof}

\begin{remark}\label{rem:schur}
Every complex irreducible representation of \(S_4\) is realizable over \(\QQ\) (Schur index \(1\)), so the hypothesis \(\End_G(W)=\QQ\) of Proposition~\ref{prop:extract} is automatically fulfilled for each of the four isotypic constituents.
\end{remark}

\begin{proposition}\label{prop:sgn-v2}
Let \(W_{\mathrm{sgn}}\) and \(W_{V_2}\) be the sign and \(V_2\)-isotypic summands of \(H^0(X,\Omega_X^1)\). Then
\[
H^0(X,\Omega_X^1)^{A_4}=W_{\mathrm{sgn}},
\qquad
H^0(X,\Omega_X^1)^{V_4}=W_{\mathrm{sgn}}\oplus W_{V_2}.
\]
Moreover, \(V_4\) acts trivially on \(V_2\), and the quotient group \(A_4/V_4\cong C_3\) acts trivially on \(W_{\mathrm{sgn}}\) and through the unique nontrivial irreducible rational representation on \(W_{V_2}\). Consequently,
\[
A_{\mathrm{sgn}}\sim_{\QQ}\Jac(X_A),
\qquad
H^0(P,\Omega_P^1)\cong W_{V_2},
\qquad
P\sim_{\QQ}A_{V_2}.
\]
\end{proposition}

\begin{proof}
From the character table computed in Theorem~\ref{thm:hodge}, the only constituent with nonzero \(A_4\)-invariants is \(\mathrm{sgn}\), contributing one invariant line per copy. Hence
\[
H^0(X,\Omega_X^1)^{A_4}=W_{\mathrm{sgn}}.
\]
For \(V_4\), the same table shows that the only constituents with nonzero invariants are \(\mathrm{sgn}\) and \(V_2\), with invariant dimensions \(1\) and \(2\), respectively. Therefore
\[
H^0(X,\Omega_X^1)^{V_4}=W_{\mathrm{sgn}}\oplus W_{V_2}.
\]

The character of \(V_2\) evaluates to \(2\) on every element of \(V_4\), so \(V_4\) acts trivially on \(W_{V_2}\). Since \(V_2\) has no \(A_4\)-invariants, the quotient \(A_4/V_4\cong C_3\) acts without fixed vectors, forcing \(W_{V_2}\) to be the unique nontrivial irreducible rational \(C_3\)-module. The sign representation is trivial on \(A_4\), so the action on \(W_{\mathrm{sgn}}\) is trivial.

Pullback identifies \(H^0(X_A,\Omega_{X_A}^1)\) with \(H^0(X,\Omega_X^1)^{A_4}=W_{\mathrm{sgn}}\). Since \(\mathrm{sgn}\) is one-dimensional, Proposition~\ref{prop:extract} gives
\[
A_{\mathrm{sgn}}\sim_{\QQ}\Jac(X_A).
\]
Likewise,
\[
H^0(X_V,\Omega_{X_V}^1)=H^0(X,\Omega_X^1)^{V_4}=W_{\mathrm{sgn}}\oplus W_{V_2}.
\]
By Corollary~\ref{cor:prym},
\[
H^0(X_V,\Omega_{X_V}^1)\cong \pi^*H^0(X_A,\Omega_{X_A}^1)\oplus H^0(P,\Omega_P^1),
\]
where \(\pi:X_V\to X_A\) is the natural cover. Comparing the two decompositions shows that
\[
H^0(P,\Omega_P^1)\cong W_{V_2}.
\]
It follows that \(P\) and \(A_{V_2}\) define the same rational Hodge substructure of \(H^1(X)\), so \(P\sim_{\QQ}A_{V_2}\).
\end{proof}

\subsection{The factorization}

\begin{theorem}\label{thm:full-factor}
Let \(E=X/S_3\) be the conductor-\(37\) elliptic quotient from Theorem~\ref{thm:elliptic}, let \(E_{\mathrm{res}}=X/D_4\) be the resolvent elliptic quotient from Proposition~\ref{prop:eres}, and let
\[
Q:=\Prym(Y/E_{\mathrm{res}}),\qquad Y=X/C_4.
\]
Then one has
\[
A_{\mathrm{sgn}}\sim_{\QQ}\Jac(X_A),\qquad
P\sim_{\QQ} E_{\mathrm{res}}^2,\qquad
A_{V_3}\sim_{\QQ} E^3,\qquad
A_{V_3^-}\sim_{\QQ} Q^3.
\]
Consequently,
\[
\Jac(X)\sim_{\QQ}\Jac(X_A)\times E_{\mathrm{res}}^2\times E^3\times Q^3.
\]
\end{theorem}

\begin{proof}
Proposition~\ref{prop:sgn-v2} identifies the sign constituent:
\[
A_{\mathrm{sgn}}\sim_{\QQ}\Jac(X_A).
\]

We identify the remaining three factors in turn.

For \(V_3\), let \(H\subset S_4\) be a point stabilizer, so that \(X/H\cong E\). The invariant table shows that \(V_3\) is the unique constituent with nonzero \(H\)-invariants, giving
\[
H^0(E,\Omega_E^1)\cong M_{V_3},
\]
where \(M_{V_3}\) is the multiplicity space of \(V_3\). Proposition~\ref{prop:extract} gives
\[
A_{V_3}\sim_{\QQ} E^3.
\]

For \(V_2\), we already know \(P\sim_{\QQ}A_{V_2}\) from Proposition~\ref{prop:sgn-v2}. The quotient \(E_{\mathrm{res}}=X/D\) satisfies
\[
H^0(E_{\mathrm{res}},\Omega_{E_{\mathrm{res}}}^1)=H^0(X,\Omega_X^1)^D,
\]
and the \(D\)-invariant column shows that this is the multiplicity space \(M_{V_2}\). Proposition~\ref{prop:extract} then gives
\[
A_{V_2}\sim_{\QQ} E_{\mathrm{res}}^2.
\]
Together with Proposition~\ref{prop:sgn-v2}, this yields
\[
P\sim_{\QQ}E_{\mathrm{res}}^2.
\]

The most delicate factor is \(V_3^-\). The degree-\(2\) map \(Y\to E_{\mathrm{res}}\) has a deck involution \(\iota\in D/C\), and the pullback decomposition gives
\[
H^0(Y,\Omega_Y^1)\cong \pi^*H^0(E_{\mathrm{res}},\Omega_{E_{\mathrm{res}}}^1)\oplus H^0(Q,\Omega_Q^1),
\]
where the first summand is the \(\iota\)-invariant part and the second is the \(\iota\)-anti-invariant part.
Now
\[
H^0(Y,\Omega_Y^1)=H^0(X,\Omega_X^1)^C.
\]
The character table gives \(\dim H^0(X,\Omega_X^1)^C=4\) (one line from \(V_2\), one from each of the three copies of \(V_3^-\)) and \(\dim H^0(X,\Omega_X^1)^D=1\) (from \(V_2\) alone). Since \(V_3^-\) has vanishing character average over \(D\), its \(C\)-invariant lines all lie in the \(\iota\)-anti-invariant eigenspace, so
\[
H^0(Q,\Omega_Q^1)\cong M_{V_3^-},
\]
and \(\dim Q=3\). Proposition~\ref{prop:extract} (applicable by Remark~\ref{rem:schur}) now gives
\[
A_{V_3^-}\sim_{\QQ} Q^3.
\]

Assembling the four identifications gives
\[
\Jac(X)\sim_{\QQ}\Jac(X_A)\times E_{\mathrm{res}}^2\times E^3\times Q^3.
\]
\end{proof}

\begin{remark}\label{rem:prym-correspondence}
The isogeny \(P\sim_{\QQ}E_{\mathrm{res}}^2\) is induced by the same rational correspondence on \(X\) that realizes the \(V_2\)-isotypic projector in \(\QQ[S_4]\). On the one hand, if \(\sigma\) generates the deck group of the \'etale cyclic triple cover \(X_V\to X_A\), then the Prym projector \(1-\sigma\) cuts out the anti-invariant part of \(\Jac(X_V)\), namely \(P\). On the other hand, the quotient map \(X\to X/D_4=E_{\mathrm{res}}\) identifies the same multiplicity space with \(H^0(E_{\mathrm{res}},\Omega_{E_{\mathrm{res}}}^1)\). The resulting pullback and norm correspondences induce homomorphisms
\[
P\longrightarrow E_{\mathrm{res}}^2,
\qquad
E_{\mathrm{res}}^2\longrightarrow P
\]
which are inverse up to finite kernel. This is the concrete \(V_2\)-block of the general functorial decomposition formalism for Jacobians of Galois covers \cite{LombardoEtAl2023}. We do not compute the minimal degree of this isogeny here.
\end{remark}

\begin{corollary}\label{cor:l-factor}
There is a corresponding factorization of first cohomology and of Hasse--Weil \(L\)-functions:
\[
H^1(X)\sim H^1(X_A)\oplus H^1(E_{\mathrm{res}})^{\oplus 2}\oplus H^1(E)^{\oplus 3}\oplus H^1(Q)^{\oplus 3},
\]
hence
\[
L(H^1(X),s)=L(H^1(X_A),s)\,L(E_{\mathrm{res}},s)^2\,L(E,s)^3\,L(H^1(Q),s)^3.
\]
\end{corollary}

\begin{proof}
This follows at once from Theorem~\ref{thm:full-factor} by passing to \(\ell\)-adic representations and taking Euler products.
\end{proof}

\begin{proposition}\label{prop:prym-end}
The Prym surface \(P=\Prym(X_V/X_A)\) carries a polarization of type \((1,3)\) and satisfies
\[
\End^0(P)\cong M_2(\QQ),
\qquad
\rho(P_{\overline{\QQ}})=3.
\]
In particular, \(P\) is not of CM type.
\end{proposition}

\begin{proof}
The Prym polarization for an \'etale cyclic triple cover of a genus-\(2\) curve has type \((1,3)\) \cite[Section~12.9]{BirkenhakeLange2004}. Since
\[
P\sim_{\QQ}E_{\mathrm{res}}^2.
\]
by Theorem~\ref{thm:full-factor}, and the \(j\)-invariant \(j(E_{\mathrm{res}})=2^{18}\cdot 7^3/(31^2\cdot 37)\) is not an algebraic integer, the curve \(E_{\mathrm{res}}\) has no complex multiplication \cite[Ch.~II]{Silverman2009}. It follows that
\[
\End^0(P)\cong \End^0(E_{\mathrm{res}}^2)\cong M_2(\QQ).
\]
Isogeny induces an isomorphism on N\'eron--Severi groups after tensoring with~\(\QQ\), so
\[
\rho(P_{\overline{\QQ}})=\rho((E_{\mathrm{res}}^2)_{\overline{\QQ}}).
\]
For the square of a non-CM elliptic curve the geometric Picard number is~\(3\), generated by the two coordinate classes and the diagonal. Since \(\End^0(P)\cong M_2(\QQ)\) is not a commutative semisimple algebra of dimension \(2\dim P\), the surface \(P\) is not of CM type.
\end{proof}

\begin{remark}\label{rem:prym-triple-cover-context}
The cover \(X_V\to X_A\) fits the cyclic triple-cover setting of Lange--Ortega \cite{LangeOrtega2011}. Their analysis of non-cyclic triple coverings of genus-\(2\) curves \cite{LangeOrtegaTriple2011} exhibits a markedly different Prym behavior, with principally polarized Prym surfaces on the non-cyclic side. In the present cyclic example the Prym polarization has type \((1,3)\), and the additional \(S_4\)-symmetry forces the stronger splitting
\[
P\sim_{\QQ}E_{\mathrm{res}}^2.
\]
\end{remark}

\section{Ray classes and weight one}

We turn to the class-field-theoretic description of the splitting field of the branch cubic. Let \(F=\QQ(\beta)\) be the cubic field generated by a root \(\beta\) of \(c(y)\), let \(L\) be the splitting field, and let
\[
K=\QQ(\sqrt{-111})
\]
be the quadratic resolvent field.

\begin{proposition}\label{prop:class-group}
The ideal class group of \(K=\QQ(\sqrt{-111})\) is cyclic of order \(8\):
\[
\Cl(K)\cong C_8.
\]
\end{proposition}

\begin{proof}
Let
\[
\tau:=\frac{1+\sqrt{-111}}{2},
\]
so \(\mathcal O_K=\ZZ[\tau]\) and \(\tau^2-\tau+28=0\). Since \(\Disc(K)=-111\), the primitive reduced positive definite binary quadratic forms of discriminant \(-111\) are
\[
(1,1,28),\ (2,\pm 1,14),\ (3,3,10),\ (4,\pm 1,7),\ (5,\pm 3,6).
\]
There are eight such forms, so \(h_K=8\).

To exhibit a generator, consider the ideal
\[
\mathfrak a:=(2,\tau-1)\subset \mathcal O_K.
\]
A direct ideal multiplication gives
\[
\mathfrak a^2=(4,\tau+3),\qquad
\mathfrak a^4=(16,\tau+11),\qquad
\mathfrak a^8=(1+3\tau).
\]
so \([\mathfrak a]^8=1\) in \(\Cl(K)\). To verify that the order is exactly \(8\), note that the norm form on \(\mathcal O_K\) is
\[
N_{K/\QQ}(x+y\tau)=x^2+xy+28y^2.
\]
If \(\mathfrak a^4\) were principal, it would be generated by an element of norm \(16\). Solving
\[
x^2+xy+28y^2=16
\]
gives only \((x,y)=(\pm 4,0)\), so the only principal ideal of norm \(16\) is \((4)\). But \(\mathfrak a^4=(16,\tau+11)\neq (4)\), so \([\mathfrak a]\) has order exactly~\(8\) and generates \(\Cl(K)\).
\end{proof}

\begin{theorem}\label{thm:conductor}
Let \(\mathfrak p_3\) be the unique prime of \(K\) above \(3\). Then:
\begin{enumerate}
\item[(i)] \(L/K\) is cyclic of degree \(3\).
\item[(ii)] The Artin conductor of \(L/K\) is
\[
\mathfrak f_{L/K}=\mathfrak p_3^2.
\]
\item[(iii)] The relative discriminant is
\[
\mathfrak d_{L/K}=\mathfrak p_3^4.
\]
\item[(iv)] The absolute discriminant of \(L\) is
\[
\Disc(L/\QQ)=3^7\cdot 37^3.
\]
\end{enumerate}
\end{theorem}

\begin{proof}
Since \(c(y)\) is irreducible with nonsquare discriminant (Proposition~\ref{prop:number-field}), the Galois group of \(L/\QQ\) is~\(S_3\), and its unique index-\(2\) subgroup fixes \(K\). Hence \(L/K\) is cyclic of degree~\(3\).

Let \(\rho\) denote the standard two-dimensional Artin representation of \(S_3\). The factorization
\[
\zeta_F(s)=\zeta(s)L(\rho,s)
\]
shows that the Artin conductor of \(\rho\) is
\[
\mathfrak f(\rho)=|\Disc(F/\QQ)|=999.
\]
On the other hand, if \(\chi\) is a nontrivial character of \(\Gal(L/K)\), then
\[
\rho\cong \operatorname{Ind}_{G_K}^{G_\QQ}\chi.
\]
The conductor formula for induction gives \cite[Ch.~VII]{Neukirch1999}
\[
\mathfrak f(\rho)=|\Disc(K/\QQ)|\,N_{K/\QQ}(\mathfrak f_\chi).
\]
Since \(|\Disc(K/\QQ)|=111\), we obtain
\[
N_{K/\QQ}(\mathfrak f_\chi)=\frac{999}{111}=9.
\]
Since \(3\) ramifies in \(K\), there is a unique prime \(\mathfrak p_3\) above \(3\), and \(\mathfrak p_3^2\) is the only ideal of norm \(9\). Therefore
\[
\mathfrak f_\chi=\mathfrak p_3^2.
\]
The two nontrivial characters of a cyclic cubic extension share the same conductor, so
\[
\mathfrak f_{L/K}=\mathfrak p_3^2
\]
and the conductor--discriminant formula \cite[Ch.~VII]{Neukirch1999} gives \(\mathfrak d_{L/K}=\mathfrak p_3^4\). Taking norms,
\[
N_{K/\QQ}(\mathfrak d_{L/K})=N(\mathfrak p_3)^4=3^4.
\]
Therefore
\[
\Disc(L/\QQ)=\Disc(K/\QQ)^3\,N(\mathfrak d_{L/K})
=111^3\cdot 3^4
=3^7\cdot 37^3.
\]
\end{proof}

\begin{remark}\label{rem:disc-distinction}
The two discriminants appearing here refer to different fields and should not be conflated. The cubic field \(F=\QQ(\beta)\) is the non-Galois cubic subfield of the \(S_3\)-extension \(L/\QQ\), and
\[
\Disc(F/\QQ)=-999=-3^3\cdot 37
\]
by Proposition~\ref{prop:number-field}. By contrast, \(L\) is the full splitting field, so its absolute discriminant is obtained from the quadratic resolvent field \(K=\QQ(\sqrt{-111})\) and the relative discriminant of the cubic extension \(L/K\):
\[
\Disc(L/\QQ)=\Disc(K/\QQ)^3\,N_{K/\QQ}(\mathfrak d_{L/K})
=111^3\cdot 3^4
=3^7\cdot 37^3.
\]
Thus items~(ii)--(iv) are mutually consistent: \(\mathfrak f_{L/K}=\mathfrak p_3^2\) implies \(\mathfrak d_{L/K}=\mathfrak p_3^4\), and after taking norms this yields the displayed discriminant of~\(L\), not the discriminant of the cubic subfield~\(F\).
\end{remark}

\begin{theorem}\label{thm:ray-class}
Let \(H_K\) be the Hilbert class field of \(K=\QQ(\sqrt{-111})\), and let \(K_{\mathfrak p_3^2}\) be the ray class field of modulus \(\mathfrak p_3^2\). Then:
\begin{enumerate}
\item[(i)] The ray class group is cyclic:
\[
\Cl_{\mathfrak p_3^2}(K)\cong C_{24}.
\]
\item[(ii)] The ray class field is the compositum
\[
K_{\mathfrak p_3^2}=H_KL,
\qquad
H_K\cap L=K.
\]
\item[(iii)] The extension \(K_{\mathfrak p_3^2}/\QQ\) is Galois with generalized dihedral group
\[
\Gal(K_{\mathfrak p_3^2}/\QQ)\cong C_{24}\rtimes C_2,
\]
where the nontrivial element of \(C_2\) acts on \(C_{24}\) by inversion.
\item[(iv)] The splitting field \(L\) of \(c(y)\) is the unique cubic subextension of \(K_{\mathfrak p_3^2}/K\).
\end{enumerate}
\end{theorem}

\begin{proof}
The standard exact sequence for the modulus \(\mathfrak p_3^2\) reads \cite[Ch.~VI]{Neukirch1999}
\[
1\longrightarrow \mathcal O_K^\times/(\mathcal O_K^\times\cap (1+\mathfrak p_3^2))
\longrightarrow (\mathcal O_K/\mathfrak p_3^2)^\times
\longrightarrow \Cl_{\mathfrak p_3^2}(K)
\longrightarrow \Cl(K)\longrightarrow 1.
\]
Since \(K\) is imaginary quadratic with discriminant \(-111\), the unit group is \(\mathcal O_K^\times=\{\pm 1\}\), and \(-1\not\equiv 1\pmod{\mathfrak p_3^2}\), so its image in \((\mathcal O_K/\mathfrak p_3^2)^\times\) has order \(2\). The residue field of \(\mathfrak p_3\) has size \(3\), so
\[
|(\mathcal O_K/\mathfrak p_3^2)^\times|
=N(\mathfrak p_3^2)\left(1-\frac{1}{N(\mathfrak p_3)}\right)
=9\left(1-\frac13\right)
=6.
\]
Equivalently, because \(3\) is ramified in \(K\), the local kernel is
\[
\frac{(\mathcal O_K/\mathfrak p_3^2)^\times}{\{\pm1\}}
\cong \frac{1+\mathfrak p_3}{1+\mathfrak p_3^2}
\cong \frac{\mathfrak p_3}{\mathfrak p_3^2}
\cong \F_3,
\]
so it is cyclic of order~\(3\). Thus the kernel of \(\Cl_{\mathfrak p_3^2}(K)\to \Cl(K)\) has order \(3\). Since \(\Cl(K)\cong C_8\) by Proposition~\ref{prop:class-group},
\[
|\Cl_{\mathfrak p_3^2}(K)|=3\cdot 8=24.
\]
Let \([\mathfrak a]\) be the generator of \(\Cl(K)\) from Proposition~\ref{prop:class-group}, and let \(u\) be any nontrivial element of the local kernel. The kernel has order coprime to \(8\), so the sequence splits and
\[
\Cl_{\mathfrak p_3^2}(K)\cong C_8\times C_3\cong C_{24},
\]
with generators \([\mathfrak a]\) and \(u\). This proves~(i).

The Hilbert class field satisfies \([H_K:K]=h_K=8\), and \(L/K\) is cyclic cubic of conductor exactly \(\mathfrak p_3^2\) (Theorem~\ref{thm:conductor}). Since \(\gcd(8,3)=1\),
\[
H_K\cap L=K.
\]
Hence
\[
[H_KL:K]=[H_K:K]\,[L:K]=24=[K_{\mathfrak p_3^2}:K],
\]
so \(K_{\mathfrak p_3^2}=H_KL\). This proves~(ii).

The modulus \(\mathfrak p_3^2\) is stable under complex conjugation, so \(K_{\mathfrak p_3^2}/\QQ\) is Galois and conjugation acts on \(\Cl_{\mathfrak p_3^2}(K)\) by inversion \cite[Ch.~9]{Cox2012}, giving~(iii).

For~(iv), \(C_{24}\) has a unique subgroup of index \(3\), and \(L/K\) corresponds to this subgroup by Theorem~\ref{thm:conductor}.
\end{proof}

\begin{remark}\label{rem:ray-class-explicit}
The proof gives a concrete presentation of the ray class group:
\[
\Cl_{\mathfrak p_3^2}(K)=\langle [\mathfrak a],u \mid [\mathfrak a]^8=u^3=1,\ [\mathfrak a]u=u[\mathfrak a]\rangle,
\]
where \(\mathfrak a=(2,\tau-1)\) and \(u\) is any nontrivial class in \((1+\mathfrak p_3)/(1+\mathfrak p_3^2)\). The companion audit file records the powers of \(\mathfrak a\) and the corresponding ray-class arithmetic used in Section~5.
\end{remark}

\begin{corollary}\label{cor:weight-one}
Let \(\rho_c\) be the standard two-dimensional Artin representation attached to the splitting field \(L/\QQ\) of \(c(y)\). Then
\[
N(\rho_c)=999,
\qquad
\det(\rho_c)=\chi_{-111},
\]
and there exists a normalized weight-\(1\) cuspidal newform
\[
f_c(q)=\sum_{n\geq 1} a_n q^n\in S_1(\Gamma_0(999),\chi_{-111})
\]
whose \(L\)-function equals \(L(\rho_c,s)\).
\end{corollary}

\begin{proof}
The Artin conductor is \(|\Disc(F/\QQ)|=999\) by Proposition~\ref{prop:number-field}, and the determinant is the quadratic character attached to \(K=\QQ(\sqrt{-111})\). Writing
\[
\rho_c\cong \operatorname{Ind}_{G_K}^{G_\QQ}\chi
\]
for a finite-order character \(\chi\) of \(\Gal(L/K)\), we see that \(\rho_c\) has finite image. Equivalently, \(\rho_c\) is the nontrivial two-dimensional summand of the permutation representation of \(G_\QQ\) on the three roots of \(c(y)\): since \(c(y)\) is irreducible with Galois group \(S_3\), this summand is irreducible. Its determinant is the sign character of \(S_3\), namely \(\chi_{-111}\), so complex conjugation acts with determinant \(-1\) and \(\rho_c\) is odd. The theorem of Deligne--Serre \cite{DeligneSerre1974WeightOne} therefore attaches a normalized weight-\(1\) cuspidal newform of level \(999\) and nebentypus \(\chi_{-111}\) whose \(L\)-function equals \(L(\rho_c,s)\).
\end{proof}

\begin{theorem}\label{thm:hecke-traces}
Let \(f_c(q)=\sum_{n\geq 1} a_n q^n\) be the normalized newform from Corollary~\ref{cor:weight-one}. Then:
\begin{enumerate}
\item[(i)] For every prime \(p\nmid 2\cdot 3\cdot 31\cdot 37\),
\[
a_p=\#\{y\in \F_p:\ c(y)=0\}-1.
\]
\item[(ii)] For every prime \(p\nmid 2\cdot 3\cdot 37\),
\[
a_p=\#\{u\in \F_p:\ u^3-9u^2+256=0\}-1.
\]
\end{enumerate}
In particular, the prime \(31\) is the unique prime away from the Artin conductor at which the maximal-order model \(u^3-9u^2+256\) still detects the Frobenius trace, while the original branch cubic \(c(y)\) does not.
\end{theorem}

\begin{proof}
For \(p\) outside the exceptional set in~(i), the reduction of \(c(y)\) modulo \(p\) is separable, so its factorization type encodes the cycle type of \(\mathrm{Frob}_p\). Let \(\rho_{\mathrm{perm}}\) be the \(3\)-dimensional permutation representation on the roots of~\(c(y)\). Then
\[
\rho_{\mathrm{perm}}\cong \mathbf{1}\oplus \rho_c,
\]
and the trace of \(\rho_{\mathrm{perm}}(\mathrm{Frob}_p)\) is the number of roots of \(c(y)\) fixed by Frobenius, i.e.\ the number of roots of \(c(y)\) in~\(\F_p\). Therefore
\[
a_p=\operatorname{tr}\rho_c(\mathrm{Frob}_p)=\operatorname{tr}\rho_{\mathrm{perm}}(\mathrm{Frob}_p)-1
=\#\mathrm{Fix}(\mathrm{Frob}_p)-1
=\#\{y\in \F_p:\ c(y)=0\}-1.
\]
The same reasoning applies to the monic model \(u^3-9u^2+256\) in~(ii). The exceptional sets differ by the single prime \(31\): the nonmonic model \(c(y)\) drops degree modulo~\(31\) (since \(256\equiv 0\)), while the maximal-order model remains separable there.
\end{proof}


\section{The genus-two Jacobian}

The genus-\(2\) factor \(\Jac(X_A)\) carries the richest arithmetic information in the decomposition. In this section we record Frobenius data at small good primes, establish the automorphic decomposition of \(L(H^1(X),s)\), obtain a Picard-number lower bound for \(\Jac(X)\), and analyze the local reduction and Galois representations of \(\Jac(X_A)\).

\subsection{Frobenius data at small good primes}

\begin{proposition}\label{thm:simple}
For the genus-two curve
\[
X_A:\quad w^2=-y(y-1)(256y^3+411y^2+165y+32),
\]
the Frobenius characteristic polynomials at the good primes \(p=5,13,17\) are
\begin{align}
\chi_5(T)&=T^4-T^3+4T^2-5T+25=(T^2+2T+5)(T^2-3T+5),\label{eq:chi5-factor}\\
\chi_{13}(T)&=T^4-2T^2+169,\label{eq:chi13}\\
\chi_{17}(T)&=T^4-3T^3+28T^2-51T+289.\label{eq:chi17}
\end{align}
In particular, \(\chi_{13}(T)\) is irreducible over~\(\QQ\). Moreover,
\[
\chi_{13}(T)=(T^2-2\sqrt{7}\,T+13)(T^2+2\sqrt{7}\,T+13)
\]
over \(\QQ(\sqrt{7})\), while
\[
\chi_{17}(T)=\left(T^2-\frac{3+\sqrt{33}}2\,T+17\right)\left(T^2-\frac{3-\sqrt{33}}2\,T+17\right)
\]
over \(\QQ(\sqrt{33})\).

More generally, for any good prime \(p\) write
\[
\chi_p(T)=T^4-s_1(p)T^3+s_2(p)T^2-s_1(p)pT+p^2.
\]
Then \(\chi_p(T)\) is reducible over~\(\QQ\) if and only if
\[
\Delta_p:=s_1(p)^2-4\bigl(s_2(p)-2p\bigr)
\]
is a square in~\(\QQ\). For
\[
p=13,17,23,29,43,47,53
\]
one finds
\[
\Delta_p=112,\ 33,\ 201,\ 73,\ 304,\ 228,\ 112,
\]
so \(\chi_p(T)\) is irreducible over~\(\QQ\) at each of these primes.
\end{proposition}

\begin{proof}
Point counting on \(X_A\) gives
\[
\#X_A(\F_5)=5,\qquad \#X_A(\F_{25})=33,
\]
\[
\#X_A(\F_{13})=14,\qquad \#X_A(\F_{169})=166,
\]
\[
\#X_A(\F_{17})=15,\qquad \#X_A(\F_{289})=337.
\]
Hence
\[
(s_1,s_2)=(1,4)\ \text{at }p=5,\qquad
(s_1,s_2)=(0,-2)\ \text{at }p=13,\qquad
(s_1,s_2)=(3,28)\ \text{at }p=17,
\]
which yields \eqref{eq:chi5-factor}, \eqref{eq:chi13}, and \eqref{eq:chi17}.

For a good prime \(p\), any factorization of
\[
\chi_p(T)=T^4-s_1T^3+s_2T^2-s_1pT+p^2
\]
over~\(\QQ\) must, by the rational-root test and Gauss's lemma, have the form
\[
\chi_p(T)=(T^2-uT+p)(T^2-vT+p),
\]
since the constant term is \(p^2\) and the reciprocal symmetry forces each quadratic factor to have constant term~\(p\). Comparing coefficients gives
\[
u+v=s_1,\qquad uv=s_2-2p.
\]
Therefore \(\chi_p(T)\) is reducible over~\(\QQ\) if and only if the quadratic equation
\[
X^2-s_1X+(s_2-2p)=0
\]
has a rational root, equivalently if and only if
\[
\Delta_p=s_1^2-4(s_2-2p)
\]
is a square in~\(\QQ\). At \(p=5\), one has \(\Delta_5=25\), which yields the displayed factorization. At \(p=13\) and \(p=17\), the values are \(\Delta_{13}=112\) and \(\Delta_{17}=33\), so both quartics are irreducible over~\(\QQ\). The same computation from the point-count data gives
\[
(s_1,s_2)=(3,-2),\ (-1,40),\ (-4,14),\ (6,46),\ (-8,94)
\]
at \(p=23,29,43,47,53\), respectively, hence the additional nonsquare values \(201,73,304,228,112\).

The displayed factorizations over \(\QQ(\sqrt{7})\) and \(\QQ(\sqrt{33})\) are obtained by solving
\[
u+v=s_1,\qquad uv=s_2-2p,
\]
for the quadratic traces \(u,v\) of Frobenius on the two degree-\(2\) factors.
\end{proof}

\begin{remark}\label{rem:frobenius-witnesses}
These repeated irreducibility witnesses do not by themselves determine the geometric endomorphism algebra of \(\Jac(X_A)\). They do, however, make accidental decompositions of \(\Jac(X_A)\) into lower-dimensional isogeny factors increasingly implausible. In particular, the extra good primes up to \(53\) reinforce the distinction between the genus-\(2\) Jacobian and the explicitly split Prym surface \(P\sim_{\QQ}E_{\mathrm{res}}^2\).
\end{remark}

\subsection{Automorphic structure}

\begin{theorem}\label{thm:automorphic}
The \(L\)-function of the first cohomology of the regular \(S_4\)-closure \(X\) admits the factorization
\begin{equation}\label{eq:automorphic-factor}
L(H^1(X),s)=L(\Jac(X_A),s)\,L(f_{1147},s)^2\,L(f_{37},s)^3\,L(H^1(Q),s)^3,
\end{equation}
where:
\begin{enumerate}
\item[(i)] \(f_{37}\in S_2(\Gamma_0(37))\) is the unique normalized newform of level \(37\), attached to the elliptic curve \(E=37\mathrm{a}1\);
\item[(ii)] \(f_{1147}\in S_2(\Gamma_0(1147))\) is the newform attached to the resolvent elliptic curve \(E_{\mathrm{res}}\), which has conductor \(1147\) and multiplicative reduction at both \(31\) and \(37\);
\item[(iii)] \(L(\Jac(X_A),s)\) is the degree-\(4\) Hasse--Weil \(L\)-function of the genus-\(2\) Jacobian \(\Jac(X_A)\);
\item[(iv)] \(Q=\Prym(Y/E_{\mathrm{res}})\) is an abelian threefold; \(H^1(Q)\) carries the \(V_3^-\)-multiplicity factor.
\end{enumerate}
\end{theorem}

\begin{proof}
Theorem~\ref{thm:full-factor} gives the Hodge-structure decomposition, and the \(L\)-function identity \eqref{eq:automorphic-factor} follows by taking Euler products. Modularity of \(L(E,s)\) and \(L(E_{\mathrm{res}},s)\) is the modularity theorem for elliptic curves over~\(\QQ\) \cite{WilesTaylor1995}. Proposition~\ref{prop:eres-conductor} identifies \(E_{\mathrm{res}}\) as an elliptic curve of conductor \(1147\). The remaining assertions are definitions of the factors appearing in the decomposition.
\end{proof}

\begin{corollary}\label{cor:potential-modularity}
The Hasse--Weil \(L\)-function \(L(\Jac(X_A),s)\) has meromorphic continuation to~\(\mathbb{C}\) and satisfies the expected functional equation.
\end{corollary}

\begin{proof}
Every abelian surface over a totally real field is potentially modular \cite{BCGP2021}; in particular \(L(\Jac(X_A),s)\) has meromorphic continuation and satisfies the expected functional equation.
\end{proof}

\begin{remark}\label{rem:structural-contrast}
The two abelian surface factors of \(\Jac(X)\) are structurally quite different. The sign factor is the Jacobian \(\Jac(X_A)\) of a concrete genus-\(2\) curve, whereas the \(V_2\)-factor \(P\sim E_{\mathrm{res}}^2\) is explicitly split and has endomorphism algebra \(M_2(\QQ)\). Thus the same branch cubic \(c(y)\) produces both a Jacobian-type abelian surface and a non-simple Prym surface.
\end{remark}

\subsection{Explicit endomorphism blocks}

\begin{proposition}\label{thm:non-isogeny}
The elliptic curves \(E\) and \(E_{\mathrm{res}}\) are not \(\QQ\)-isogenous. Moreover,
\[
\End^0(\Jac(X)_{\overline{\QQ}})
\supseteq
\End^0(\Jac(X_A)_{\overline{\QQ}})
\times M_2(\QQ)\times M_3(\QQ)\times \End^0(Q^3_{\overline{\QQ}}).
\]
\end{proposition}

\begin{proof}
For the first assertion, \(\QQ\)-isogenous elliptic curves have the same conductor. Since
\[
\mathrm{cond}(E)=37\neq 1147=\mathrm{cond}(E_{\mathrm{res}}),
\]
the curves \(E\) and \(E_{\mathrm{res}}\) are not \(\QQ\)-isogenous.

The displayed containment is immediate from the \(\QQ\)-isogeny
\[
\Jac(X)\sim_{\QQ}\Jac(X_A)\times E_{\mathrm{res}}^2\times E^3\times Q^3
\]
of Theorem~\ref{thm:full-factor}. The matrix blocks \(M_2(\QQ)\) and \(M_3(\QQ)\) come from the standard endomorphisms of the powers \(E_{\mathrm{res}}^2\) and \(E^3\).
\end{proof}

\begin{corollary}\label{cor:picard}
The geometric Picard number of \(\Jac(X)\) satisfies
\[
\rho(\Jac(X)_{\overline{\QQ}})\geq 16.
\]
\end{corollary}

\begin{proof}
The isogeny decomposition of Theorem~\ref{thm:full-factor} yields
\[
\rho(\Jac(X)_{\overline{\QQ}})\geq \rho(\Jac(X_A)_{\overline{\QQ}})+\rho(E_{\mathrm{res}}^2{}_{\overline{\QQ}})+\rho(E^3_{\overline{\QQ}})+\rho(Q^3_{\overline{\QQ}}).
\]
Now \(\rho(\Jac(X_A)_{\overline{\QQ}})\geq 1\). For the non-CM elliptic curves \(E_{\mathrm{res}}\) and \(E\), the Rosati-symmetric parts of \(M_2(\QQ)\) and \(M_3(\QQ)\) have dimensions \(3\) and \(6\), so
\[
\rho(E_{\mathrm{res}}^2{}_{\overline{\QQ}})=3,\qquad \rho(E^3_{\overline{\QQ}})=6.
\]
Finally, \(Q^3\) carries the obvious matrix endomorphisms from the three coordinate factors, so its Rosati-symmetric endomorphisms contain the \(6\)-dimensional space of symmetric \(3\times 3\) rational matrices. Hence \(\rho(Q^3_{\overline{\QQ}})\geq 6\), and therefore
\[
\rho(\Jac(X)_{\overline{\QQ}})\geq 1+3+6+6=16.
\]
\end{proof}

\subsection{The Prym threefold}

\begin{proposition}\label{prop:prym-threefold}
The abelian threefold \(Q=\Prym(Y/E_{\mathrm{res}})\) satisfies:
\begin{enumerate}
\item[(i)] \(\Jac(Y)\sim_{\QQ}E_{\mathrm{res}}\times Q\).
\item[(ii)] The \(\ell\)-adic Galois representation on the first cohomology of \(Q\) is
\[
V_\ell(Q)\cong M_{V_3^-}\otimes \QQ_\ell,
\]
where \(M_{V_3^-}\) is the multiplicity space of the anti-standard representation \(V_3^-=V_3\otimes\mathrm{sgn}\) of~\(S_4\).
\item[(iii)] The \(L\)-function of \(Q\) is a degree-\(6\) Euler product. The factorization~\eqref{eq:automorphic-factor} gives
\[
L(H^1(Q),s)^3=\frac{L(H^1(X),s)}{L(\Jac(X_A),s)\,L(E_{\mathrm{res}},s)^2\,L(E,s)^3},
\]
so \(L(H^1(Q),s)\) is determined by the \(L\)-functions of the remaining factors and the genus-\(16\) curve \(X\).
\end{enumerate}
\end{proposition}

\begin{proof}
Part~(i) is the standard Prym decomposition for the degree-\(2\) cover \(Y\to E_{\mathrm{res}}\) \cite{Mumford1974,LangeOrtega2011}.

For~(ii), Theorem~\ref{thm:full-factor} identifies \(H^0(Q,\Omega_Q^1)\) with \(M_{V_3^-}\); the same identification holds at the \(\ell\)-adic level. Part~(iii) follows by dividing \eqref{eq:automorphic-factor} by the three known factors and comparing Euler products at unramified primes.
\end{proof}

\begin{proposition}\label{prop:prym-traces}
For every prime \(p\nmid 2\cdot 3\cdot 31\cdot 37\), the Frobenius trace \(a_1(Q,p)\) is determined by
\[
a_1(Q,p)=\frac{a_1(X,p)-a_1(\Jac(X_A),p)-2\,a_1(E_{\mathrm{res}},p)-3\,a_1(E,p)}{3},
\]
where \(\#X(\F_p)=24\,N_{\mathrm{split}}(p)+12(2+c(p))+6\), with \(N_{\mathrm{split}}(p)\) the number of \(y\in\F_p\setminus\{0,1\}\) for which \(c(y)\neq 0\) and \(\Pi(\lambda,y)\) splits completely, and \(c(p)=\#\{y\in\F_p:c(y)=0\}\). The first fifteen values are:
\begin{center}
% Auto-generated by scripts/generate_Q_traces_table.py
\begin{tabular}{r|rrrr}
$p$ & $a_1(\Jac(X_A))$ & $a_1(E_{\mathrm{res}})$ & $a_1(E)$ & $a_1(Q)$ \\ \hline
$5$ & $1$ & $-2$ & $-2$ & $-5$ \\
$7$ & $3$ & $-5$ & $-1$ & $-4$ \\
$11$ & $6$ & $-3$ & $-5$ & $-5$ \\
$13$ & $0$ & $-2$ & $-2$ & $-6$ \\
$17$ & $3$ & $0$ & $0$ & $-5$ \\
$19$ & $2$ & $-6$ & $0$ & $-4$ \\
$23$ & $3$ & $0$ & $2$ & $-5$ \\
$29$ & $-1$ & $-4$ & $6$ & $-3$ \\
$41$ & $10$ & $-5$ & $-9$ & $1$ \\
$43$ & $-4$ & $6$ & $2$ & $-12$ \\
$47$ & $6$ & $-9$ & $-9$ & $-9$ \\
$53$ & $-8$ & $-5$ & $1$ & $-7$ \\
$59$ & $-3$ & $-6$ & $8$ & $7$ \\
$61$ & $4$ & $8$ & $-8$ & $-8$ \\
$67$ & $1$ & $-4$ & $8$ & $-9$ \\
\end{tabular}

\end{center}
All computed values satisfy the Weil bound \(|a_1(Q,p)|\leq 6\sqrt{p}\) and the divisibility \(3\mid (a_1(X,p)-a_1(\Jac(X_A),p)-2\,a_1(E_{\mathrm{res}},p)-3\,a_1(E,p))\), providing a numerical consistency check for the isogeny decomposition of Theorem~\ref{thm:full-factor}.
\end{proposition}

\begin{proof}
For the regular \(S_4\)-closure \(\pi\colon X\to\PP^1\), the fiber above an unramified rational point \(y\in\PP^1(\F_p)\) consists of \(|S_4|=24\) geometric points on which Frobenius acts by left multiplication. Since left multiplication by a non-identity element of \(S_4\) has no fixed points, the fiber contributes \(24\) rational points if and only if the Frobenius is trivial, i.e.\ \(\Pi(\lambda,y)\) splits completely; otherwise the fiber contributes~\(0\). The branch locus consists of \(y=0\), \(y=1\) (transposition inertia, \(|S_4/\langle\sigma\rangle|=12\) points each), the roots of \(c(y)\) (transposition inertia, \(12\) points each), and \(y=\infty\) (\(4\)-cycle inertia, \(|S_4/\langle\tau\rangle|=6\) points).

By Theorem~\ref{thm:automorphic},
\[
3\,a_1(Q,p)=a_1(X,p)-a_1(\Jac(X_A),p)-2\,a_1(E_{\mathrm{res}},p)-3\,a_1(E,p).
\]
The traces \(a_1(E,p)\) and \(a_1(E_{\mathrm{res}},p)\) are obtained by point counting on the explicit Weierstrass models of \(E\) and \(E_{\mathrm{res}}\), while \(a_1(\Jac(X_A),p)\) is computed from the hyperelliptic model \(w^2=-y(y-1)c(y)\). Divisibility by~\(3\) and the Weil bound are verified for all primes \(p<120\) with \(p\nmid 2\cdot 3\cdot 31\cdot 37\).
\end{proof}

\subsection{Conductor and reduction}

\begin{proposition}\label{prop:disc-quintic}
The discriminant of the hyperelliptic model \(f(y)=-y(y-1)c(y)\) of \(X_A\) is
\[
\Disc(f)=-2^{20}\cdot 3^{15}\cdot 31^2\cdot 37.
\]
The set of primes of bad reduction for \(X_A\) is contained in \(\{2,3,31,37\}\).
\end{proposition}

\begin{proof}
The roots of \(f(y)=-y(y-1)c(y)\) are \(0,1,\beta_1,\beta_2,\beta_3\), where \(\beta_i\) are the roots of \(c(y)\). For leading coefficient \(a_5=-256\),
\[
\Disc(f)=a_5^{8}\prod_{i<j}(r_i-r_j)^2.
\]
The product of pairwise differences factors as
\begin{align*}
\prod_{i<j}(r_i-r_j)^2
&=(0-1)^2\cdot\prod_i(0-\beta_i)^2\cdot\prod_i(1-\beta_i)^2\cdot\prod_{i<j}(\beta_i-\beta_j)^2.
\end{align*}
By Vieta's formulas for \(c(y)=256\prod_i(y-\beta_i)\):
\[
\prod_i\beta_i=-\frac{32}{256}=-\frac{1}{8},\qquad
\prod_i(1-\beta_i)=\frac{c(1)}{256}=\frac{864}{256}=\frac{27}{8}.
\]
Using \(\Disc(c)=256^4\prod_{i<j}(\beta_i-\beta_j)^2=-3^9\cdot 31^2\cdot 37\) (Proposition~\ref{prop:number-field}),
\[
\prod_{i<j}(r_i-r_j)^2=1\cdot\frac{1}{64}\cdot\frac{729}{64}\cdot\frac{-3^9\cdot 31^2\cdot 37}{256^4}
=\frac{-3^{15}\cdot 31^2\cdot 37}{2^{44}}.
\]
Therefore \(\Disc(f)=(-256)^8\cdot(-3^{15}\cdot 31^2\cdot 37)/2^{44}=2^{64}\cdot(-3^{15}\cdot 31^2\cdot 37)/2^{44}=-2^{20}\cdot 3^{15}\cdot 31^2\cdot 37\).
\end{proof}

\begin{theorem}\label{thm:reduction-types}
The reduction types of the genus-\(2\) curve \(X_A\) at the bad primes are as follows:
\begin{enumerate}
\item[(i)] At \(p=31\): the polynomial \(f(y)\bmod 31\) has exactly one double root (at \(y\equiv 12\)), giving semistable reduction with one ordinary double point. The conductor exponent is \(f_{31}=1\).
\item[(ii)] At \(p=37\): the polynomial \(f(y)\bmod 37\) has exactly one double root (at \(y\equiv 6\)), giving semistable reduction with one ordinary double point. The conductor exponent is \(f_{37}=1\).
\item[(iii)] At \(p=3\): the branch cubic collapses to \(c(y)\equiv (y-1)^3\pmod{3}\), so
\[
f(y)\equiv 2y(y-1)^4\pmod{3},
\]
and the special fiber has a non-split tacnode at \((1,0)\) of delta invariant~\(2\): the two local branches are conjugate over~\(\F_9\) and meet to second order. Moreover, the three non-rational branch points form a single depth-\(3/2\) cluster above \(y=1\), so \(X_A\) does not satisfy the semistability criterion of \cite[Definition~1.8 and Theorem~1.9]{DDMM2023}; in particular \(X_A/\QQ_3\) is not semistable.
\item[(iv)] The conductor of \(\Jac(X_A)\) has the form
\[
N_A=2^{f_2}\cdot 3^{f_3}\cdot 31\cdot 37,\qquad f_2\geq 1,
\]
and Theorem~\ref{thm:conductor-exact} identifies the wild part at \(3\) as \(\delta_3=1\).
\end{enumerate}
\end{theorem}

\begin{proof}
For~(i), one computes \(c(y)\equiv 8(y-6)(y-12)^2\pmod{31}\), so
\[
f(y)\equiv -8y(y-1)(y-6)(y-12)^2\pmod{31},
\]
with simple roots \(0,1,6\) and double root~\(12\). The double root produces a single ordinary node on the special fiber. For a tame prime with a single node, Liu's analysis \cite{Liu2002} gives \(f_p=1\).

For~(ii), modulo \(37\) the branch cubic factors as \(c(y)\equiv -3(y-6)^2(y-14)\pmod{37}\), giving
\[
f(y)\equiv 3y(y-1)(y-6)^2(y-14)\pmod{37}.
\]
and the same argument gives \(f_{37}=1\).

For~(iii), the branch cubic satisfies \(c(y)\equiv (y-1)^3\pmod{3}\), so the special fiber of \(w^2=f(y)\) over~\(\F_3\) is \(w^2=2y(y-1)^4\). The singularity at \((1,0)\) is a non-split tacnode: substituting \(y=1+t\) gives \(w^2=2t^4(1+t)\), whose two local branches \(w=\pm\sqrt{2}\,t^2\sqrt{1+t}\) are conjugate over~\(\F_9\) (since \(\sqrt{2}\notin\F_3\)) and meet to second order, giving delta invariant~\(2\). Proposition~\ref{prop:cluster-pictures}(i) shows that the three cubic branch points form a depth-\(3/2\) cluster requiring a totally ramified cubic extension of~\(\QQ_3\) to separate. Since the semistability criterion of \cite[Definition~1.8]{DDMM2023} requires the root field \(K(\mathcal R)\) to have ramification degree at most~\(2\), Theorem~1.9 of loc.\ cit.\ implies that \(X_A/\QQ_3\) is not semistable.

For~(iv), modulo \(2\) the polynomial \(f\) reduces to \(y^2(y+1)^2\), which is a perfect square, so \(f_2\geq 1\). The wild part at \(3\) is identified in Theorem~\ref{thm:conductor-exact}.
\end{proof}

\begin{proposition}\label{prop:node-type}
The nodal singularities of the special fibers of \(X_A\) at the two semistable primes have opposite splitting types:
\begin{enumerate}
\item[(i)] At \(p=31\), the node is split: the tangent directions at the double root \(y\equiv 12\) are rational over~\(\F_{31}\). The normalization of the nodal special fiber is the elliptic curve \(v^2=-8y(y-1)(y-6)\) over~\(\F_{31}\).
\item[(ii)] At \(p=37\), the node is non-split: the tangent directions at the double root \(y\equiv 6\) are conjugate over~\(\F_{37^2}\). The normalization is the elliptic curve \(v^2=3y(y-1)(y-14)\) over~\(\F_{37}\).
\end{enumerate}
In both cases, the Jacobian \(\Jac(X_A)\) acquires toric rank~\(1\) at each prime, with the abelian part of the N\'eron model equal to the Jacobian of the normalization.
\end{proposition}

\begin{proof}
Write \(f(y)=(y-a)^2g(y)\) at the double root \(y\equiv a\pmod{p}\). The tangent cone of \(w^2=f(y)\) at the node is \(v^2=g(a)\) (after \(v=w/(y-a)\)), which splits over~\(\F_p\) if and only if \(g(a)\in(\F_p^\times)^2\).

At \(p=31\), \(a=12\): we have \(g(y)=-8y(y-1)(y-6)\) and
\[
g(12)\equiv -8\cdot 12\cdot 11\cdot 6\equiv 19\equiv 9^2\pmod{31},
\]
so the node is split.

At \(p=37\), \(a=6\): here \(g(y)=3y(y-1)(y-14)\) and
\[
g(6)\equiv 3\cdot 6\cdot 5\cdot(-8)\equiv 20\pmod{37}.
\]
The Legendre symbol \(\left(\frac{20}{37}\right)=\left(\frac{5}{37}\right)=-1\), so \(g(6)\) is a quadratic non-residue and the node is non-split.

In both cases \(v^2=g(y)\) is an elliptic curve over~\(\F_p\), and the single node gives toric rank~\(1\) in the N\'eron model \cite{Liu2002}.
\end{proof}

\begin{proposition}\label{prop:cluster-pictures}
The \(p\)-adic cluster pictures of the roots of \(f(y)=-y(y-1)c(y)\) at the two wild primes are as follows.
\begin{enumerate}
\item[(i)] At \(p=3\): the Newton polygon of the translated polynomial \(c(1+3t)/27=256t^3+393t^2+195t+32\) centered at the collapsed root \(y=1\) shows that the three roots \(\beta_i\) of \(c(y)\) satisfy \(v_3(\beta_i-1)=1\) and \(v_3(\beta_i-\beta_j)=3/2\) for all \(i\neq j\). The cluster picture is
\[
\{0\}\;\bigg|\;\underbrace{\{1\}\;\bigg|\;\underbrace{\{\beta_1,\beta_2,\beta_3\}}_{\text{depth }3/2}}_{\text{depth }1},
\]
where the roots \(\beta_i\) are conjugate over~\(\QQ_3\) and require a totally ramified cubic extension to separate. Since this extension has degree~\(3=p\), it contributes a nontrivial Swan conductor \(\delta_3>0\).
\item[(ii)] At \(p=2\): the Newton polygon of \(c(y)=256y^3+411y^2+165y+32\) has vertices at \((0,v_2(32))\!=\!(0,5)\), \((1,v_2(165))\!=\!(1,0)\), \((2,v_2(411))\!=\!(2,0)\), \((3,v_2(256))\!=\!(3,8)\), giving slopes \(-5,0,8\). (With the convention that the Newton polygon slope between \((i,v_i)\) and \((i+1,v_{i+1})\) is \(v_{i+1}-v_i\), each root \(\beta\) satisfies \(v_2(\beta)=-\text{slope}\) of the corresponding segment.) Each segment has length~\(1\), so by Hensel's lemma all three roots lie in~\(\QQ_2\). Denote these roots \(\beta_1,\beta_2,\beta_3\). Since \(v_2(c(1))=5\) and \(v_2(c'(1))=0\), one has \(v_2(\beta_2-1)=5\). The cluster picture is
\[
\underbrace{\{0,\beta_1\}}_{\text{depth }5}\;\bigg|\;\underbrace{\{1,\beta_2\}}_{\text{depth }5}\;\bigg|\;\{\beta_3\},
\]
with two balanced clusters of depth~\(5\) and one isolated root of valuation \(-8\).
\end{enumerate}
\end{proposition}

\begin{proof}
For~(i), the substitution \(y=1+3t\) gives \(c(1+3t)=27(256t^3+393t^2+195t+32)=27h(t)\). A further substitution \(t=1+u\) yields
\[
h(1+u)=256u^3+1161u^2+1749u+876,
\]
with \(3\)-adic valuations \(v_3(876)=1\), \(v_3(1749)=1\), \(v_3(1161)=3\), \(v_3(256)=0\). The Newton polygon has a single segment of slope \(-1/3\) from \((0,1)\) to \((3,0)\), so all three roots satisfy \(v_3(t_i-1)=1/3\). In particular \(v_3(\beta_i-1)=1+v_3(t_i)=1\), and
\[
v_3(\beta_i-\beta_j)=1+v_3(t_i-t_j).
\]
The pairwise distances follow from \(\Disc(h)=-3^3\cdot 31^2\cdot 37\), giving
\[
\textstyle\sum_{i<j}v_3(t_i-t_j)=\tfrac12\bigl(v_3(\Disc(h))-4\cdot v_3(256)\bigr)=\tfrac32,
\]
and by symmetry each \(v_3(t_i-t_j)=1/2\), whence \(v_3(\beta_i-\beta_j)=3/2\).

For~(ii), the \(2\)-adic Newton polygon of \(c(y)\) has vertices \((0,5)\), \((1,0)\), \((2,0)\), \((3,8)\), giving slopes \(-5\), \(0\), \(8\). The root \(\beta_1\) with \(v_2(\beta_1)=5\) clusters with~\(0\); the root \(\beta_2\) with \(v_2(\beta_2)=0\) satisfies \(v_2(\beta_2-1)=5\) by Newton's method applied to \(c(1+h)=0\), since \(v_2(c(1))=5\) and \(v_2(c'(1))=0\); the root \(\beta_3\) with \(v_2(\beta_3)=-8\) is isolated.
\end{proof}

\begin{proposition}\label{lem:local-three-prep}
Let \(F_w=\QQ_3(\beta_1)\), where \(\beta_1\) is any root of the branch cubic \(c(y)\). Then:
\begin{enumerate}
\item[(i)] \(F_w/\QQ_3\) is totally ramified of degree~\(3\);
\item[(ii)] the local discriminant exponent is \(v_3(\Delta_{F_w/\QQ_3})=3\);
\item[(iii)] the finite branch roots of \(f(y)=-y(y-1)c(y)\) form exactly three \(G_{\QQ_3}\)-orbits, namely
\[
\{0\},\qquad \{1\},\qquad \{\beta_1,\beta_2,\beta_3\}.
\]
\end{enumerate}
\end{proposition}

\begin{proof}
By Proposition~\ref{prop:cluster-pictures}(i), the roots \(\beta_i\) are conjugate over~\(\QQ_3\) and satisfy \(v_3(\beta_i-1)=1\) and \(v_3(\beta_i-\beta_j)=3/2\). In particular the denominator~\(2\) in the pairwise distances forces a nontrivial ramified extension, and the same proposition shows that a totally ramified cubic extension is required to separate the three roots. Thus \([F_w:\QQ_3]=3\) and \(F_w/\QQ_3\) is totally ramified, proving~(i). The orbit statement~(iii) is the same cluster-theoretic description, together with the rationality of the roots \(0\) and~\(1\).

For~(ii), Proposition~\ref{prop:number-field} gives \(\Disc(F/\QQ)=-3^3\cdot 37\). Since \(F_w/\QQ_3\) is totally ramified of degree~\(3\), the prime \(3\) has a unique extension to~\(F\) and residue degree~\(1\). Therefore the exponent of \(3\) in the global field discriminant equals the local discriminant exponent:
\[
v_3(\Disc(F/\QQ))=f_w\,v_3(\Delta_{F_w/\QQ_3})=1\cdot v_3(\Delta_{F_w/\QQ_3}),
\]
so \(v_3(\Delta_{F_w/\QQ_3})=3\).
\end{proof}

\begin{theorem}\label{thm:conductor-exact}
The wild part of the conductor of \(\Jac(X_A)\) at \(p=3\) is exactly
\[
\delta_3=1.
\]
\end{theorem}

\begin{proof}
Let
\[
R=\{0,1,\beta_1,\beta_2,\beta_3\}
\]
be the set of finite branch roots of the odd-degree model
\[
w^2=-256\prod_{r\in R}(y-r).
\]
The branch polynomial has odd degree \(5\), the residue characteristic is \(p=3\neq 2\), and the roots are distinct in \(\overline{\QQ}_3\) because \(\Disc(f)\neq 0\). Therefore the conductor formula of Dokchitser, Dokchitser, Maistret, and Morgan (Math. Ann. 385 (2023), Theorem~11.3) applies exactly in this setting: for an odd-degree hyperelliptic curve, the wild conductor is expressed as a sum over the \(G_{\QQ_3}\)-orbits on the finite branch set \(R\), with no extra contribution from the point at infinity. Since the scalar factor \(-1\) is a \(3\)-adic unit, it does not affect the extension fields \(\QQ_3(r)\) or the conductor contribution of any root. Thus
\[
\delta_3=\sum_{r\in R//G_{\QQ_3}}
\Bigl(v_3(\Delta_{\QQ_3(r)/\QQ_3})-[\QQ_3(r):\QQ_3]+f(\QQ_3(r)/\QQ_3)\Bigr),
\]
where \(R//G_{\QQ_3}\) denotes a set of \(G_{\QQ_3}\)-orbit representatives and \(f(\QQ_3(r)/\QQ_3)\) is the residue degree.

By Proposition~\ref{lem:local-three-prep}, the \(G_{\QQ_3}\)-orbits on \(R\) are \(\{0\}\), \(\{1\}\), and \(\{\beta_1,\beta_2,\beta_3\}\). The rational roots \(0\) and~\(1\) contribute
\[
v_3(\Delta_{\QQ_3/\QQ_3})-[\QQ_3:\QQ_3]+f(\QQ_3/\QQ_3)=0-1+1=0.
\]
Equivalently, the cluster picture contributes two rational singleton orbits and one cubic orbit centered \(3\)-adically at~\(1\), with pairwise distances \(v_3(\beta_i-\beta_j)=3/2\). All of the wild contribution is therefore concentrated in that single cubic orbit.
For the cubic orbit, Proposition~\ref{lem:local-three-prep} gives \([F_w:\QQ_3]=3\), \(f(F_w/\QQ_3)=1\), and \(v_3(\Delta_{F_w/\QQ_3})=3\).
Therefore the cubic orbit contributes
\[
v_3(\Delta_{F_w/\QQ_3})-[F_w:\QQ_3]+f(F_w/\QQ_3)=3-3+1=1,
\]
so the full sum is
\[
\delta_3=(0-1+1)+(0-1+1)+(3-3+1)=1.
\]
\end{proof}

\begin{remark}\label{rem:f2-obstruction}
The conductor exponent \(f_2\) remains undetermined. At \(p=2\), the special fiber of the naive hyperelliptic model is \(w^2\equiv y^2(y-1)^2\pmod{2}\), a non-reduced curve, so the minimal regular model over~\(\ZZ_2\) requires a nontrivial sequence of blowups whose combinatorics depend on the full \(2\)-adic cluster picture. Since all three roots of \(c(y)\) lie in~\(\QQ_2\) (Proposition~\ref{prop:cluster-pictures}(ii)), the stable reduction is in principle computable, but the two balanced clusters at depth~\(5\) make the computation substantially more involved than the wild-part calculation at~\(p=3\).
\end{remark}

\subsection{Torsion and Galois representations}

\begin{theorem}\label{thm:two-torsion}
The \(2\)-torsion field of \(\Jac(X_A)\) coincides with the splitting field of the branch cubic:
\[
\QQ(\Jac(X_A)[2])=L.
\]
The mod-\(2\) Galois representation
\[
\bar\rho_2\colon G_{\QQ}\longrightarrow \mathrm{Sp}(4,\F_2)\cong S_6
\]
has image isomorphic to \(S_3\), of order~\(6\) inside \(|\mathrm{Sp}(4,\F_2)|=720\).
\end{theorem}

\begin{proof}
For a genus-\(2\) curve \(C:w^2=f(y)\) with \(\deg f=5\), the nonzero elements of \(\Jac(C)[2]\) are represented by divisor classes of the form
\[
[(r_i,0)-(r_j,0)],
\]
where \(r_i,r_j\) are Weierstrass points; equivalently, the \(2\)-torsion field is the splitting field of~\(f\) \cite[Ch.~2]{CasselsFlynn1996}. In our case
\[
f(y)=-y(y-1)c(y)
\]
has Weierstrass points above \(0\), \(1\), and the three roots \(\beta_1,\beta_2,\beta_3\) of~\(c(y)\). Since \(0\) and~\(1\) are rational, adjoining all \(2\)-torsion points amounts to adjoining the cubic branch points \(\beta_i\). Hence
\[
\QQ(\Jac(X_A)[2])=\QQ(\beta_1,\beta_2,\beta_3)=L.
\]
Let \(W_0=\langle D_0,D_1\rangle\), where \(D_0=[(0,0)-\infty]\) and \(D_1=[(1,0)-\infty]\). Then \(W_0\) is fixed pointwise by \(G_{\QQ}\). In the quotient \(\Jac(X_A)[2]/W_0\), the classes \(\bar D_{\beta_1},\bar D_{\beta_2}\) form an \(\F_2\)-basis and satisfy \(\bar D_{\beta_3}=\bar D_{\beta_1}+\bar D_{\beta_2}\). Therefore a \(3\)-cycle \(\sigma=(123)\) and a transposition \(\tau=(12)\) act by
\[
\sigma\mapsto
\begin{pmatrix}
0&1\\
1&1
\end{pmatrix},
\qquad
\tau\mapsto
\begin{pmatrix}
0&1\\
1&0
\end{pmatrix}
\]
on \(\Jac(X_A)[2]/W_0\). These matrices generate \(\mathrm{GL}(2,\F_2)\cong S_3\), so the mod-\(2\) image is exactly \(S_3\). Since the Weil pairing is Galois invariant, this image embeds into \(\mathrm{Sp}(4,\F_2)\) as a subgroup of order~\(6\).
\end{proof}

\begin{corollary}\label{cor:two-torsion-classfield}
Theorem~\ref{thm:ray-class}(iv) identifies \(L\) as the unique cubic subextension of \(K_{\mathfrak p_3^2}/K\). The \(2\)-torsion of \(\Jac(X_A)\) is therefore controlled by the ray class group \(\Cl_{\mathfrak p_3^2}(K)\cong C_{24}\), and its absolute discriminant is \(\Disc(L/\QQ)=3^7\cdot 37^3\) (Theorem~\ref{thm:conductor}(iv)). The same extension~\(L/\QQ\) also determines the weight-\(1\) Artin representation of Corollary~\ref{cor:weight-one}.
\end{corollary}

\begin{proof}
Combine Theorem~\ref{thm:two-torsion} with Theorem~\ref{thm:ray-class}(iv).
\end{proof}

\subsection{The \texorpdfstring{$\F_2[S_3]$}{F2[S3]}-module structure of the \texorpdfstring{$2$}{2}-torsion}

Theorem~\ref{thm:two-torsion} identifies the \(2\)-torsion field as the splitting field~\(L\) of~\(c(y)\) and shows that the mod-\(2\) image is~\(S_3\). We now refine this to a complete description of the module structure and the resulting embedding.

\begin{theorem}\label{thm:mod2-structure}
As an \(\F_2[\Gal(L/\QQ)]\)-module, the \(2\)-torsion admits a decomposition
\[
\Jac(X_A)[2]\cong W_0\oplus W_1,
\]
where \(W_0=\langle D_0,D_1\rangle\cong \F_2^2\) is the trivial module spanned by the rational Weierstrass divisor classes and \(W_1\cong \mathrm{St}_2\) is the standard two-dimensional representation of~\(S_3\cong \mathrm{GL}(2,\F_2)\). Both \(W_0\) and \(W_1\) are nondegenerate (hyperbolic) subspaces for the Weil pairing, and the decomposition is a symplectic orthogonal sum:
\[
\bigl(\Jac(X_A)[2],e_2\bigr)\cong (W_0,e_2|_{W_0})\perp (W_1,e_2|_{W_1}).
\]
Consequently, the mod-\(2\) Galois image lies in the subgroup
\[
\bar\rho_2(G_{\QQ})\cong S_3\hookrightarrow \{1\}\times \mathrm{Sp}(W_1)\subset \mathrm{Sp}(W_0)\times \mathrm{Sp}(W_1)\subset \mathrm{Sp}(4,\F_2),
\]
acting trivially on the rational hyperbolic plane \(W_0\) and faithfully on the irrational hyperbolic plane~\(W_1\).
\end{theorem}

\begin{proof}
For the genus-\(2\) curve \(w^2=f(y)\) with \(\deg f=5\), the \(2\)-torsion \(\Jac(X_A)[2]\cong (\ZZ/2\ZZ)^4\) is generated by the divisor classes \(D_i=[(r_i,0)-\infty]\) for the five roots \(r_i\) of~\(f\), subject to \(\sum D_i=0\) \cite{CasselsFlynn1996}. The roots are \(0,1,\beta_1,\beta_2,\beta_3\). Taking \(\{D_0,D_1,D_{\beta_1},D_{\beta_2}\}\) as a basis, the Galois group \(S_3=\Gal(L/\QQ)\) fixes \(D_0\) and \(D_1\), and permutes \(\{D_{\beta_1},D_{\beta_2},D_{\beta_3}\}\) with \(D_{\beta_3}=D_0+D_1+D_{\beta_1}+D_{\beta_2}\).

For a genus-\(2\) Jacobian in a Weierstrass basis, the Weil pairing on \(2\)-torsion satisfies \(e_2(D_i,D_j)=-1\) for \(i\neq j\) \cite{CasselsFlynn1996}. Hence the Gram matrix in additive \(\F_2\)-notation with respect to \(\{D_0,D_1,D_{\beta_1},D_{\beta_2}\}\) is
\[
M=\begin{pmatrix}0&1&1&1\\1&0&1&1\\1&1&0&1\\1&1&1&0\end{pmatrix}.
\]
Consider the symplectic basis \(e_1=D_0,\;f_1=D_1,\;e_2=D_{\beta_1}+D_{\beta_2},\;f_2=D_0+D_1+D_{\beta_2}\). A direct computation gives
\[
\langle e_1,f_1\rangle=1,\quad \langle e_2,f_2\rangle=1,\quad \langle e_i,e_j\rangle=\langle f_i,f_j\rangle=\langle e_i,f_j\rangle=0\quad (i\neq j),
\]
exhibiting the standard symplectic form \(J_2\perp J_2\). The plane \(W_0=\langle e_1,f_1\rangle\) is fixed by~\(S_3\), and its symplectic complement \(W_1=\langle e_2,f_2\rangle\) carries the standard representation (the quotient action on \(\{D_{\beta_1},D_{\beta_2},D_{\beta_3}\}\) modulo~\(W_0\) is the standard \(2\)-dimensional representation of \(S_3\cong \mathrm{GL}(2,\F_2)\)). Both summands are hyperbolic, and the decomposition is \(S_3\)-equivariant.

The subgroup of \(\mathrm{Sp}(4,\F_2)\) preserving the orthogonal splitting \(W_0\perp W_1\) is \(\mathrm{Sp}(W_0)\times\mathrm{Sp}(W_1)\cong S_3\times S_3\). Since Theorem~\ref{thm:two-torsion} identifies the quotient action on \(W_1\) with the standard matrices above and \(S_3\) acts trivially on \(W_0\), the image is \(\{1\}\times S_3\).
\end{proof}

\begin{remark}\label{rem:hyperbolic-splitting}
The \(S_3\)-invariant hyperbolic splitting \(W_0\perp W_1\) is the unique such decomposition: the only \(2\)-dimensional \(S_3\)-invariant nondegenerate subspace is \(W_0=\Jac(X_A)[2]^{G_{\QQ}}\), and its symplectic complement~\(W_1\) is the unique \(S_3\)-stable complement. This rigidity contrasts with the situation for curves with fewer rational Weierstrass points, where the fixed space may be trivial and the Galois representation cannot be reduced in this manner.
\end{remark}

\subsection{Component groups at the semistable primes}

\begin{proposition}\label{prop:tamagawa}
The component groups and Tamagawa numbers of \(\Jac(X_A)\) are as follows:
\begin{enumerate}
\item[(i)] At \(p=31\) (split node): the component group is \(\Phi_{31}\cong \ZZ/2\ZZ\), and \(c_{31}(\Jac(X_A))=2\).
\item[(ii)] At \(p=37\) (non-split node): the component group is trivial, and \(c_{37}(\Jac(X_A))=1\).
\end{enumerate}
\end{proposition}

\begin{proof}
At each semistable prime, the stable model of \(X_A\) has a single ordinary double point (Theorem~\ref{thm:reduction-types}). Blowing up that node once produces the minimal regular model: the special fiber consists of the strict transform \(C_0\) of the normalization, which is an elliptic curve, together with one exceptional divisor \(E_0\cong \PP^1\). The exceptional divisor records the two tangent directions at the node.

For a one-nodal irreducible semistable curve, Raynaud's description of the N\'eron model identifies the identity component of the Jacobian with an extension of \(\Jac(C_0)\) by a one-dimensional torus determined by those tangent directions; see \cite[Chapter~10]{BoschLutkebohmertRaynaud1990} and \cite[Chapter~10]{Liu2002}. If the node is split, the two tangent directions are rational and the torus is \(\mathbb{G}_m\); the monodromy pairing has matrix \([2]\), so the component group is
\[
\Phi_p\cong \operatorname{coker}\bigl([2]\colon \ZZ\to \ZZ\bigr)\cong \ZZ/2\ZZ.
\]
If the node is non-split, the tangent directions are exchanged by Frobenius and the toric part is the anisotropic norm-one torus \(R^1_{\F_{p^2}/\F_p}\mathbb{G}_m\), whose N\'eron model is connected. In that case the component group is trivial.

By Proposition~\ref{prop:node-type}, the node at~\(31\) is split, giving \(\Phi_{31}\cong \ZZ/2\ZZ\) and \(c_{31}=|\Phi_{31}(\F_{31})|=2\). The node at~\(37\) is non-split, giving \(\Phi_{37}=0\) and \(c_{37}=1\).
\end{proof}

\subsection{Numerical verification}

\begin{proposition}\label{prop:verification}
The following table records the point counts \(\#X_A(\F_p)\) and \(\#E_{\mathrm{res}}(\F_p)\) at the first twelve primes of good reduction for~\(X_A\), together with the Frobenius traces and the number of roots of \(c(y)\) modulo~\(p\):
\[
\begin{array}{c|cccc|c}
p & \#X_A(\F_p) & a_1(X_A) & \#E_{\mathrm{res}}(\F_p) & a_1(E_{\mathrm{res}}) & \#\{y\in\F_p:c(y)=0\}\\ \hline
5 & 5 & 1 & 8 & -2 & 0 \\
7 & 5 & 3 & 13 & -5 & 0 \\
11 & 6 & 6 & 15 & -3 & 1 \\
13 & 14 & 0 & 16 & -2 & 1 \\
17 & 15 & 3 & 18 & 0 & 0 \\
19 & 18 & 2 & 26 & -6 & 1 \\
23 & 21 & 3 & 24 & 0 & 0 \\
29 & 31 & -1 & 34 & -4 & 0 \\
41 & 32 & 10 & 47 & -5 & 1 \\
43 & 48 & -4 & 38 & 6 & 1 \\
47 & 42 & 6 & 57 & -9 & 1 \\
53 & 62 & -8 & 59 & -5 & 1 \\
\end{array}
\]
The characteristic polynomial of Frobenius \(\chi_{13}(T)=T^4-2T^2+169\) is irreducible over~\(\QQ\), in agreement with Proposition~\ref{thm:simple}.
\end{proposition}

\begin{proof}
The point counts are computed by evaluating the Legendre symbol \((f(y)/p)\) for each \(y\in\F_p\), together with the point at infinity. The characteristic polynomial is determined by \(a_1=p+1-\#X_A(\F_p)\) and \(a_2=(a_1^2-p^2-1+\#X_A(\F_{p^2}))/2\). Irreducibility of \(\chi_{13}(T)=T^4-2T^2+169\) is shown in the proof of Proposition~\ref{thm:simple}.
\end{proof}

\begin{remark}\label{rem:verification-hecke}
The Hecke eigenvalues \(a_p(f_c)=\#\{y\in\F_p:c(y)=0\}-1\) of the weight-\(1\) form, computed from the last column of the table, match the values predicted by Theorem~\ref{thm:hecke-traces} at all primes listed. The values cycle through \(\{-1,0,2\}\), corresponding to the Frobenius cycle types \((123)\), \((12)\), and~\(1\) in~\(S_3\).
\end{remark}

\begin{remark}\label{rem:reproducibility-scripts}
For a structured inventory of the companion scripts and generated files used in the numerical checks of Sections~5 and~6, see the reproducibility note recorded after Section~7.
\end{remark}


\section{Further directions}

\begin{enumerate}
\item[(1)] Determine \(\End^0(\Jac(X_A)_{\overline{\QQ}})\) and the odd \(\ell\)-adic monodromy of \(\Jac(X_A)\).

\item[(2)] Construct an explicit model for \(Y=X/C_4\) and determine the resulting geometry and \(L\)-function of the Prym threefold \(Q=\Prym(Y/E_{\mathrm{res}})\).

\item[(3)] Proposition~\ref{prop:cluster-pictures} determines the \(2\)-adic and \(3\)-adic cluster pictures, and Theorem~\ref{thm:conductor-exact} identifies the exact wild contribution \(\delta_3=1\) at~\(3\). Determining the full conductor at~\(3\) and the exponent \(f_2\) still requires a full stable-model computation or an independent automorphic input.
\end{enumerate}

\subsection*{Reproducibility}
The computational claims used in Sections~5 and~6 are accompanied by scripts in the source tree of the manuscript. The main files are:
\begin{itemize}
\item \texttt{point\_count.py} in \texttt{artifacts/}, which reproduces the point counts of \(X_A\) over \(\F_p\) and \(\F_{p^2}\);
\item \texttt{conductor\_exact.py} in \texttt{artifacts/}, which checks the local conductor calculation at \(p=3\);
\item \texttt{exp\_branch\_cubic\_rayclass\_modform\_audit.py} in \texttt{scripts/}, which reproduces the class-group, ray-class, and Hecke-trace checks of Section~5;
\item \texttt{exp\_genus2\_jacobian\_audit.py} in \texttt{scripts/}, which audits the Frobenius polynomials of \(\Jac(X_A)\);
\item \texttt{generate\_Q\_traces\_table.py} in \texttt{scripts/}, which generates the table for~\(Q\).
\end{itemize}
The companion note \texttt{REPRODUCIBILITY.md} records the role of each script and the generated outputs used in the manuscript. Database information from the LMFDB is used only as a consistency check and not as an input to any proof.

\subsection*{Acknowledgements}
Computations were performed with SageMath and Python. The LMFDB~\cite{LMFDB} was used as a consistency check for elliptic-curve metadata.

\subsection*{Declaration of generative AI and AI-assisted technologies in the manuscript preparation process}
During the preparation of this work the authors used OpenAI ChatGPT in order to improve the language, readability, and presentation of the manuscript, without generating the underlying mathematical ideas, arguments, or results. After using this tool, the authors reviewed and edited the content as needed and take full responsibility for the content of the published article.

\bibliographystyle{unsrt}
\bibliography{references}

\end{document}
